\documentclass[11pt,a4paper]{article}
\usepackage[utf8]{inputenc}
\usepackage[T1]{fontenc}
\usepackage{lmodern}
\usepackage[margin=1in]{geometry}
\usepackage{amsmath,amssymb,amsthm}
\usepackage{booktabs}
\usepackage{longtable}
\usepackage{array}
\usepackage{microtype}
\usepackage[hidelinks]{hyperref}

\newtheorem{theorem}{Theorem}
\newtheorem{proposition}[theorem]{Proposition}
\newtheorem{lemma}[theorem]{Lemma}
\newtheorem{corollary}[theorem]{Corollary}
\theoremstyle{definition}
\newtheorem{definition}[theorem]{Definition}
\theoremstyle{remark}
\newtheorem{remark}[theorem]{Remark}

\newcommand{\lQ}{\lambda_{Q}}
\newcommand{\lR}{\lambda_{R}}
\newcommand{\lA}{\lambda_{A}}
\newcommand{\PsiL}{\Psi_{L}}
\newcommand{\src}[1]{{\footnotesize\texttt{#1}}}
\newcommand{\st}[1]{{\small\textsf{[#1]}}}

\title{\bfseries Certified Weil Positivity in Two Windows:\\
Spectral Bounds, Conditioning, and Reproducible Certificates}

\author{Elias De Jes\'us\\[2pt]
{\small Independent Researcher}\\[1pt]
{\small ORCID \href{https://orcid.org/0009-0007-0190-9143}{0009-0007-0190-9143}}\\[1pt]
{\small \texttt{dejesuselias10@gmail.com}}}

\date{November 2026 --- Technical Note (preprint), version 2.0.0\\[4pt]
{\small First archived release. Continues the internal draft series 1.0.x--1.2.0;
see Section~\ref{sec:history}.}}

\begin{document}
\maketitle

\begin{center}
\fbox{\begin{minipage}{0.93\textwidth}
\small
\textbf{Scope.} This note certifies two-sided bounds on the spectral infimum of
the Weil explicit-formula quadratic form restricted to a compact window, at two
explicit windows, together with per-sector conditioning results and one
eigenvalue count. Every decisive inequality is produced in Arb ball arithmetic
from exact rational inputs and is replayed by a release gate shipped with the
package (Section~\ref{sec:repro}).

\smallskip
\textbf{Not claimed.} No arithmetic estimate, no zero-free region, no verified
range of zeros, no statement about the Riemann hypothesis, and no result at any
window larger than those stated. No priority is claimed over the literature
cited in Section~\ref{sec:prior}. The certified lower bounds do not extend above
their stated windows.

\smallskip
\textbf{Excluded.} An earlier compact time-domain evaluator omitted the
truncation error of its outer panelled Gauss--Legendre rule. That route is
\emph{excluded} from this release and every quantity depending on it remains
uncertified, retained only as a labelled measurement
(Section~\ref{sec:status}). Nothing in the results below uses it.
\end{minipage}}
\end{center}

\begin{abstract}
\noindent
Let $Q$ be the Weil explicit-formula quadratic form restricted to test functions
supported in a compact window $[-L,L]$, in one fixed normalization, and let
$\lQ(L)$ be its spectral infimum. This note reports computer-assisted two-sided
bounds at two windows,
\[
   1.43300048827856\times 10^{-17} \;\le\; \lQ(4/5) \;\le\; 1.674\times 10^{-17},
\]
\[
   8.27343868254787\times 10^{-18} \;\le\; \lQ(\widetilde L) \;\le\; 9.475\times 10^{-18},
   \qquad \widetilde L = \tfrac{80471895621705018}{10^{17}} < \tfrac{\log 5}{2}.
\]
Each lower bound comes from a frequency-cutoff comparison at $T = 400$, reduced
to a $440$-dimensional Legendre block whose positive definiteness is verified by
a preconditioned congruence in Arb ball arithmetic, together with a
super-exponentially small complement estimate and a strict Schur-type coupling
test; both parity sectors are certified at each window. Each upper bound is the
Rayleigh quotient of the \emph{original} form at one explicit even rational test
function, evaluated in closed form through the digamma, Lerch $\Phi$ and Hurwitz
$\zeta$ functions --- a route with no quadrature, hence no truncation term to
bound. Because the form restricts consistently to smaller windows, $\lQ$ is
non-increasing in $L$; the first lower bound therefore holds throughout
$0 < L \le 4/5$, the second throughout $0 < L \le \widetilde L$, and the second
upper bound for every $L \ge \widetilde L$. Pairing the four endpoints gives a
corollary on the \emph{true} minima,
$0.494 \le \lQ(\widetilde L)/\lQ(4/5) \le 0.662$. At $L = 4/5$ we also certify,
per parity sector, that removing four explicit rational directions from the even
sector raises the infimum to at least $0.6175$, and that exactly four even
eigenvalues lie below $10^{-2}$ --- at most four from the conditioning
certificate, at least four from a certified four-dimensional test space for $Q$
itself. In the odd sector only ``at most three'' is certified. A proved
normalization dictionary identifies this form with the window form of
Andrews--Connes--La~Gatta under $\lambda = e^{L}$, verified numerically to about
thirty digits. All constants are published as midpoint--radius enclosures with a
release gate that replays the chain from the packaged files. This is a
computation at two windows: it establishes no arithmetic estimate, nothing about
the location of any zeta zero, and nothing about the Riemann hypothesis.
\end{abstract}

\section{Introduction}

\subsection{The object}

Weil's explicit formula pairs a test function against the zeros of $\zeta$, the
pole at $s=1$, the archimedean place and the primes. When the test function is
required to be supported in a fixed compact window, the resulting quadratic form
becomes a concrete object: a bounded-below Hermitian form on $L^2(-L,L)$ whose
positivity \emph{for every} $L$ is equivalent to the Riemann
hypothesis~\cite{Weil,Bombieri}, and whose positivity at one fixed $L$ is not. Finite-window Weil
positivity in this sense is the object studied in the spectral programs of
Connes and van Suijlekom~\cite{CvS}, Connes and Consani~\cite{CC21}, and
Connes, Consani and Moscovici~\cite{CCM}, and quantitatively in Zhu~\cite{Zhu}, which supplies the
finite-reduction strategy that the present computation follows.

The quantity computed here is the spectral infimum
\[
   \lQ(L) \;=\; \inf_{f} \frac{Q_{\zeta,L}(f)}{\|f\|_2^2}
\]
of that form at $L = 4/5$, in the normalization fixed in
Section~\ref{sec:setup}. Two things make a two-sided statement possible at this
window and difficult at larger ones. An \emph{upper} bound needs a single
admissible test function and is unconditional. A \emph{lower} bound needs the
whole window, hence a reduction of an infinite-dimensional positivity question
to a finite matrix plus certified tails; the cost of that reduction grows very
fast in $L$, which is the content of Theorem 1.4 of~\cite{Zhu}.

\subsection{What this note reports}

\begin{enumerate}
\item A certified lower bound $Q(f) \ge 1.4330\times10^{-17}\|f\|_2^2$ at
  $L = 4/5$, in both parity sectors and hence for arbitrary complex tests
  (Section~\ref{sec:lower}).
\item A certified upper bound $\lQ(4/5) \le 1.674\times10^{-17}$ from one
  explicit even rational test function, with the \emph{original} form evaluated
  in closed form (Section~\ref{sec:upper}). A time-domain candidate is reported
  but withdrawn pending a quadrature bound (Section~\ref{sec:status}).
\item A certified parity separation: the odd sector satisfies
  $\lambda_{Q,\mathrm{odd}} \ge 1.3856\times10^{-14}$, which is strictly above the
  upper bound, so the even infimum determines $\lQ$
  (Section~\ref{sec:parity}).
\item Cutoff diagnostics at $T \in \{200,400,800\}$ that measure what the
  comparison discards on fixed test functions, and a certified ceiling on what
  a larger cutoff can deliver (Section~\ref{sec:cutoff}).
\end{enumerate}

\subsection{Relation to prior work, and what is not claimed}
\label{sec:prior}

The finite-to-infinite reduction strategy --- a scalar gate on the exterior of a
frequency cutoff, a Legendre/Bessel representation of the interior, and
super-exponentially small tails --- is Theorem 1.1 of~\cite{Zhu}. Zhu describes
supplementary source code and certificates in~\cite[\S17]{Zhu}; these were not
obtained or executed in the present work, so his matrix step is neither re-run
nor verified here. The reduction was instead rederived in the normalization of
the present repository, and every constant below was computed from scratch. The bracket $8.9\times10^{-18} \le \lambda_{\min}(0.8) \le
2.27\times10^{-17}$ stated in the abstract of~\cite{Zhu} was used only as an
external comparison point and never as an input, an expected answer, or a
tolerance. The parity structure is likewise not new, and at this window it is
not merely the decoupling that is already published: \cite[Thm.~6.2]{Zhu}
certifies the \emph{separation} at support $1.6$ --- that is, at $L = 4/5$ --- and
does so with a far stronger odd-sector constant than the one obtained here,
concluding in addition that the ground state is simple and even. What
Section~\ref{sec:parity} contributes is an independent re-derivation of the
ordering from this pipeline's own certified data, nothing structural.

The computation continues the implementation developed alongside our
finite-window sensitivity study~\cite{Sensitivity}, whose repository supplies the
zeta-window normalization of Section~\ref{sec:setup}, the geometric-side
assembly, the pole conventions, and the closed-form archimedean evaluator used
as the validation reference in Section~\ref{sec:upper}(a). The even/odd pole
signs used here are the ones checked entrywise against that assembly in our
real-zero parity study~\cite{ParityNote}. Neither the $\delta^2$--$L^3$
sensitivity estimate of~\cite{Sensitivity} nor the planted-real-zero analysis
of~\cite{ParityNote} is an input to the lower bound proved here: the first
concerns a counterfactual perturbation of the zero set, and the second a
configuration that is nowhere assumed to exist. The finite-to-infinite reduction
follows~\cite{Zhu}; what is contributed here is its execution in certified
arithmetic, together with the resulting bracket, the numerical parity separation
of Section~\ref{sec:parity}, and the cutoff diagnostics of
Section~\ref{sec:cutoff}.

What follows is deliberately narrow, and the following are \emph{not} claimed.
No arithmetic estimate is obtained: nothing here improves any bound on primes,
on $\Lambda(n)$, on $\psi(x)$, or on a zero-free region. No statement is made
about the location of any zero of $\zeta$, and no range of zeros is verified.
No new positivity framework is introduced; the framework is that of the cited
sources, and the contribution is a certified computation inside it. The
computation is at one window, $L = 4/5$, and gives nothing at any larger support
--- there is no propagation theorem here, and a certificate at one $L$ does not
provide one. Nothing here bears on the Riemann hypothesis. This is a technical
note, not peer-reviewed research, and it is not claimed that anyone else uses or
has used this work.

\subsection{The same window, compared}
\label{sec:compared}

Because~\cite{Zhu} certifies at support $1.6$, its results and the present ones
describe the same object in the same normalization (its equations (2)--(3) are
\eqref{eq:Q}--\eqref{eq:Psi} below), so they can be put side by side. Doing so is
the fairest summary of what this note adds and what it does not.

\begin{center}
\begin{tabular}{lrr}
\toprule
quantity at $L = 4/5$ & \cite{Zhu}, certified & this note \\
\midrule
whole-window infimum
  & $[8.9\times10^{-18},\, 2.27\times10^{-17}]$
  & $[1.4330\times10^{-17},\, 1.674\times10^{-17}]$ \\
odd-sector lower bound
  & $8.206\times10^{-15}$
  & $1.3856\times10^{-14}$ \\
second even min--max value
  & $\ge 2.085\times10^{-12}$
  & not certified here \\
simplicity of the ground state
  & \cite[Thm.~6.2]{Zhu}
  & not independently established here \\
\bottomrule
\end{tabular}
\end{center}

Two rows improve on the published result. The width of the whole-window bracket
falls from $1.38\times10^{-17}$ to $2.4013\times10^{-18}$: approximately
$82.6\%$ narrower, a reduction by a factor of $5.747$, with the present bracket
strictly nested inside the published one, both endpoints moving inward. The
odd-sector constant is larger by a factor $1.68$. Neither is a new method: both
are sharper constants obtained inside the framework of~\cite{Zhu}, from the same
kind of certificate. In the remaining two rows~\cite{Zhu} is stronger, and the
simplicity and evenness of the ground state at this window --- which this note
does not establish --- are his.

As a consistency check --- not as an endpoint, and not as a replay of his
certificate --- \cite{Zhu} also reports a numerical reference estimate
$\lambda_{\min}(Q|_{\text{window}}) = 1.656\times10^{-17}$ for this window, from
his own exact time-domain assembly. That value lies inside the bracket certified
here. It is a non-certified reference value in~\cite{Zhu} and is used here only
in that capacity.

\section{Setup, conventions, and the statement}
\label{sec:setup}

\subsection{Normalization}

Fix $L > 0$. For $f \in L^2(-L,L)$, extended by zero, write
\[
   F(t) \;=\; \int_{-L}^{L} f(u)\, e^{itu}\, du ,
   \qquad\text{so}\qquad
   \int_{\mathbb{R}} |F(t)|^2\, dt \;=\; 2\pi \|f\|_2^2 .
\]
The window Weil form used throughout is the Hermitian form
\begin{equation}
\label{eq:Q}
   Q(f) \;=\; 2\operatorname{Re}\!\big(F(i/2)\,\overline{F(-i/2)}\big)
     \;+\; \frac{1}{2\pi}\int_{\mathbb{R}} |F(t)|^2\, \PsiL(t)\, dt ,
\end{equation}
\begin{equation}
\label{eq:Psi}
   \PsiL(t) \;=\; H(t) - P_L(t), \qquad
   H(t) \;=\; \operatorname{Re}\psi\!\left(\tfrac14 + \tfrac{it}{2}\right) - \log\pi ,
   \qquad
   P_L(t) \;=\; \sum_{\log n < 2L} \frac{2\Lambda(n)}{\sqrt n}\,\cos(t \log n) .
\end{equation}
Here $\psi$ is the digamma function and $\Lambda$ is the von Mangoldt function.
Writing $C = \int f\cosh(u/2)\,du$ and $S = \int f\sinh(u/2)\,du$, so that
$F(\mp i/2) = C \pm S$, the pole term of~\eqref{eq:Q} is
\begin{equation}
\label{eq:pole}
   2\operatorname{Re}\!\big(F(i/2)\,\overline{F(-i/2)}\big)
   \;=\; 2\big(|C|^2 - |S|^2\big) ,
\end{equation}
which is the form that must be used for complex $f$: it is additive over the
real and imaginary parts, whereas $2F(i/2)^2$ is not --- replacing $f$ by $if$
reverses its sign. The computation never needs~\eqref{eq:pole} in that
generality. It is carried out on \emph{real} test functions of definite parity,
where $C$ survives on even $f$ and $S$ on odd $f$, and~\eqref{eq:pole}
specializes to
\[
   2 s\, F(i/2)^2 , \qquad s = +1 \text{ (even)}, \qquad s = -1 \text{ (odd)} ,
\]
the two signs recording that the pole helps the even sector and hurts the odd
one. Arbitrary complex $f$ is then reached through the decomposition of
Section~\ref{sec:parity-complex}, not through~\eqref{eq:pole} directly.

The two sector signs are not asserted here: in the
parity study~\cite[\S3(c)]{ParityNote} the $\delta = 1/2$ real-pair perturbation
is compared entrywise with this repository's own pole term
$2\,\mathrm{pole\_sign}(\varepsilon)P_jP_k$ and agrees in all nine tested cells to
$6.22\times10^{-61}$, with $\mathrm{pole\_sign} = -1$ on the odd sector, while
the opposite sign convention is wrong by $O(1)$ to $O(10)$;
and~\cite[Remark 6]{ParityNote} identifies that $\delta = 1/2$ case as the parity
splitting of~\cite[Lemma 6.1]{Zhu}. The prime sum in~\eqref{eq:Psi} is
\emph{finite and exact}: the autocorrelation of a function supported in
$[-L,L]$ is supported in $[-2L,2L]$, so only $n$ with $\log n < 2L$ contribute
and there is no truncation. At $L = 4/5$ the active prime powers are
$\{2,3,4\}$.

This is the normalization of \src{src/zeta\_window.py} in the repository
accompanying this note, documented in~\cite{Sensitivity}, where the same object
is written $T_R = A - 2M$ with $A = \mathrm{Pole} + \mathrm{Arch} - (\log\pi)I$
and $2M = \mathrm{Prime}$. It is an imported convention, not derived here, and
the grouping of $-\log\pi$ is part of it. It was cross-checked against the repository's own
independent closed form for the constant test function
(Section~\ref{sec:upper}, item~(a)).

\subsection{The admissible class}

Since $H$ is even, continuous, and increasing on $t > 0$
(Lemma~\ref{lem:gate}), and $|P_L| \le A_L := \sum_{\log n < 2L}
2\Lambda(n)/\sqrt n < \infty$, the function $\PsiL$ is bounded below on
$\mathbb{R}$ and bounded on compacts. Hence for \emph{every} $f \in L^2(-L,L)$
the negative part of the integral in~\eqref{eq:Q} is finite and $Q(f)$ is
well defined in $(-\infty,+\infty]$. We therefore take the admissible class to
be all of $L^2(-L,L)$, with the convention that a lower bound
$Q(f) \ge \mu\|f\|_2^2$ is satisfied trivially where $Q(f) = +\infty$. This is a
choice of convention made in this note; it strengthens rather than weakens the
statement, and it removes any need for a density or form-domain argument.

\subsection{Three spectral quantities, kept apart}

Confusing the following three numbers is the easiest way to overstate a result
of this kind, so they are named separately and used separately throughout.

\begin{definition}
\label{def:three}
With $T > 0$ a frequency cutoff, $\beta$ a constant with $\PsiL(t) \ge \beta$
for $|t| \ge T$, and $V_N$ a finite-dimensional subspace of $L^2(-L,L)$:
\begin{align*}
  \lQ      &= \inf\{\,Q(f)/\|f\|_2^2 \,\}
             && \text{over the admissible class,}\\
  \lR(T)   &= \inf\{\,R_{T,\beta}(f)/\|f\|_2^2 \,\}
             && \text{over the same class, } R_{T,\beta} \text{ as in \eqref{eq:R},}\\
  \lA(T,N,\varepsilon) &= \lambda_{\min}\big(\text{compression of } R_{T,\beta} \text{ to } V_N\big)
             && \text{in the parity sector } \varepsilon .
\end{align*}
\end{definition}

\noindent Because $Q \ge R_{T,\beta}$ pointwise, and because an infimum over a
subspace is at least the infimum over the whole space,
\[
   \lR(T) \;\le\; \lQ , \qquad \lR(T) \;\le\; \lA(T,N,\varepsilon) ,
\]
and there is \textbf{no} inequality in either direction between $\lA$ and $\lQ$.
In particular a numerical least eigenvalue of a finite block is not an upper
bound for $\lQ$, and only the evaluation of $Q$ itself on an explicit test
function is.

\subsection{The theorem}

\begin{theorem}
\label{thm:main}
Let $Q$ be the form~\eqref{eq:Q}--\eqref{eq:Psi} on $L^2(-L,L)$ and
$\lQ(L) = \inf\{Q(f)/\|f\|_2^2\}$. Put
\[
   \widetilde L \;=\; \frac{80471895621705018}{10^{17}} ,
   \qquad\text{so}\qquad \tfrac45 \;<\; \widetilde L \;<\; \frac{\log 5}{2} .
\]
\begin{enumerate}
\item[(i)] \emph{(Lower bound at $L = 4/5$.)} For every
  $f \in L^2(-4/5,4/5)$,
  \[
     Q(f) \;\ge\; 1.43300048827856\times 10^{-17}\, \|f\|_2^2 ,
  \]
  in both parity sectors separately and hence for arbitrary complex $f$.
\item[(ii)] \emph{(Upper bound at $L = 4/5$.)} There is an explicit even test
  function $f_{160}$, a combination of the $160$ even cosine modes
  $\cos(k\pi u/L)$, $k = 0,\dots,159$, with rational coefficients of
  denominator $10^{180}$, such that
  \[
     \lQ(4/5) \;\le\; \frac{Q(f_{160})}{\|f_{160}\|_2^2}
     \;\le\; U_0 \;=\; 1.67312553407068186807726837736\times10^{-17}
     \;<\; 1.674\times10^{-17}.
  \]
\item[(iii)] \emph{(Lower bound at $\widetilde L$.)} For every
  $f \in L^2(-\widetilde L,\widetilde L)$,
  \[
     Q(f) \;\ge\; 8.27343868254787\times 10^{-18}\, \|f\|_2^2 ,
  \]
  again in both parity sectors separately.
\item[(iv)] \emph{(Upper bound at $\widetilde L$.)} There is an explicit even
  test function $\widetilde f_{160}$ of the same shape such that
  \[
     \lQ(\widetilde L) \;\le\; U_- \;=\;
     9.47476951083049032849709\times10^{-18} \;<\; 9.475\times10^{-18}.
  \]
\item[(v)] \emph{(Parity separation at $L = 4/5$.)} For every odd
  $f \in L^2(-4/5,4/5)$,
  $Q(f) \ge 1.38569999999861\times10^{-14}\|f\|_2^2$. Since $f_{160}$ is even
  and $1.3856\times10^{-14} > U_0$, the even infimum equals $\lQ(4/5)$.
\end{enumerate}
\end{theorem}

Parts (i), (iii) and (v) are computer-assisted: they rest on interval arithmetic
in the Arb midpoint--radius model~\cite{Arb}, including all analytic error
terms, on exact rational input data, and on a release gate that recomputes every
decisive implication from the packaged files. Parts (ii) and (iv) each need one
test function and are unconditional given the normalization.

\begin{remark}
The shortened decimals above are directed: an upper bound is quoted rounded
\emph{up} and a lower bound rounded \emph{down}. In particular
$1.674\times10^{-17}$ is a valid short form for $U_0$ while
$1.6731\times10^{-17}$ is \emph{not}, being below $U_0$; and $1.4330\times
10^{-17}$ is a valid short form for the constant in (i). Section~\ref{sec:repro}
records the check that enforces this on every printed bound.
\end{remark}

\begin{remark}
Theorem~\ref{thm:main}(v) identifies which sector carries the infimum. It does
\emph{not} assert that the infimum is attained, that a minimizer is unique, or
that the least eigenvalue is simple.
\end{remark}

\begin{proposition}[Monotonicity and coverage]
\label{prop:mono}
If $0 < L' \le L$ then every $f$ supported in $[-L',L']$ is admissible for the
window $[-L,L]$ and $Q$ takes the same value on it, so
$\lQ(L) \le \lQ(L')$: the infimum is non-increasing in $L$. Consequently
Theorem~\ref{thm:main} gives
\[
   \lQ(L) \;\ge\; 1.43300048827856\times10^{-17}
   \quad (0 < L \le \tfrac45) , \qquad
   \lQ(L) \;\ge\; 8.27343868254787\times10^{-18}
   \quad (0 < L \le \widetilde L) ,
\]
\[
   \lQ(L) \;\le\; 9.475\times10^{-18} \quad (L \ge \widetilde L) ,
   \text{ in particular at } L = \tfrac{\log 5}{2} .
\]
No positive lower bound at $L = (\log 5)/2$ follows from these certificates:
monotonicity runs the other way.
\end{proposition}

\begin{corollary}[Change between the two minima]
\label{cor:ratio}
With $\mu_0, U_0$ and $\mu_-, U_-$ the certified lower and upper endpoints at
$4/5$ and $\widetilde L$,
\[
   \frac{\mu_-}{U_0} \;\le\; \frac{\lQ(\widetilde L)}{\lQ(4/5)}
   \;\le\; \frac{U_-}{\mu_0} ,
   \qquad\text{hence}\qquad
   0.494 \;\le\; \frac{\lQ(\widetilde L)}{\lQ(4/5)} \;\le\; 0.662 ,
\]
so the spectral infimum falls by between $33.8\%$ and $50.6\%$ between the two
windows.
\end{corollary}

\begin{proof}
Both brackets are two-sided, and $\lQ > 0$ at both windows by
Theorem~\ref{thm:main}(i),(iii), so the quotient is bounded by pairing the lower
endpoint at each window with the upper endpoint at the other. Taking the four
packaged exact decimals $\mu_0 = 1.43300048827856\times10^{-17}$,
$U_0 = 1.674\times10^{-17}$, $\mu_- = 8.27343868254787\times10^{-18}$,
$U_- = 9.475\times10^{-18}$ gives
$\mu_-/U_0 = 0.4942317\ldots$ and $U_-/\mu_0 = 0.6612000\ldots$, and the
displayed bounds are those rounded outward. The divisions are exact rational
arithmetic.
\end{proof}

\begin{remark}
Corollary~\ref{cor:ratio} is a consequence of the two brackets and nothing else:
no prolate quantity, no parameter transfer, and no numerical eigenvalue from the
literature enters it. It is \emph{not} an independent analytic scaling theorem,
and it does not establish any asymptotic law. See
Section~\ref{sec:dict} for what it does and does not say about
Hypothesis~3.4 of~\cite{ACLG}.
\end{remark}

\section{The certified lower bound}
\label{sec:lower}

The lower bound is assembled in five steps. Steps 1--5 are proved here in the
normalization of Section~\ref{sec:setup}; the strategy is that
of~\cite[Thm.~1.1]{Zhu}.

\subsection{Step 1: the exterior gate}

\begin{lemma}
\label{lem:gate}
$t \mapsto \operatorname{Re}\psi(\tfrac14 + \tfrac{it}{2})$ is strictly
increasing on $t > 0$.
\end{lemma}

\begin{proof}
From the digamma series,
\[
   \operatorname{Re}\psi\!\left(\tfrac14+\tfrac{it}{2}\right)
   = -\gamma + \sum_{n \ge 0}
     \left[\frac{1}{n+1} - \frac{n + 1/4}{(n+1/4)^2 + t^2/4}\right],
\]
and each summand is strictly increasing in $t > 0$. Termwise differentiation is
justified by local uniform convergence.
\end{proof}

Consequently, with $A_L = \sum_{\log n < 2L} 2\Lambda(n)/\sqrt n$ and
$\beta := H(T) - A_L$, we get $\PsiL(t) \ge \beta$ for all $|t| \ge T$, because
$P_L \le A_L$ pointwise. At $L = 4/5$,
\[
   A_L = 2.941973525223620455503909 \pm 3.91\times10^{-25},
\]
and the gate values at the three cutoffs used in this note are
\begin{center}
\begin{tabular}{lll}
\toprule
$T$ & $\beta = H(T) - A_L$ & status \\
\midrule
$200$ & $0.5184657332438467199873536 \pm 2.62\times10^{-26}$ & certified \\
$400$ & $1.211613695058064549471502 \pm 1.89\times10^{-26}$ & certified \\
$800$ & $1.904761070930776888621206 \pm 1.68\times10^{-25}$ & certified \\
\bottomrule
\end{tabular}
\end{center}
with the increments $\beta_{400} - \beta_{200} = 0.6931479618142178294841484 \pm
7.24\times10^{-27}$, numerically close to $\log 2$ as the asymptotics of $H$
suggest. All three are \emph{positive}, so the set
$\{|t| \ge T : \PsiL(t) < \beta\}$ is empty and no occupation or measure
estimate on a bad exterior set is needed anywhere below.

\subsection{Step 2: the comparison form}

Because $\int_{|t| \ge T} (\PsiL - \beta)|F|^2 \ge 0$, equation~\eqref{eq:Q}
gives $Q(f) \ge R_{T,\beta}(f)$ with
\begin{equation}
\label{eq:R}
   R_{T,\beta}(f) \;=\; \beta \|f\|_2^2 \;+\; 2 s F(i/2)^2
     \;+\; \frac{1}{2\pi}\int_{-T}^{T} \big(\PsiL(t) - \beta\big) |F(t)|^2\, dt .
\end{equation}
The discarded quantity is exactly
\begin{equation}
\label{eq:e}
   e_T(f) \;:=\; \frac{Q(f) - R_{T,\beta}(f)}{\|f\|_2^2}
   \;=\; \frac{1}{2\pi\|f\|_2^2}\int_{|t| > T} \big(\PsiL(t) - \beta\big)|F(t)|^2\, dt
   \;\ge\; 0 ,
\end{equation}
which is measured in Section~\ref{sec:cutoff}.

\begin{proposition}[monotonicity in $\beta$]
\label{prop:beta}
By Plancherel, $R_{T,\beta}(f) = 2sF(i/2)^2 + \frac{1}{2\pi}\int_{|t|\le T}
\PsiL |F|^2 + \frac{\beta}{2\pi}\int_{|t| > T}|F|^2$, so for
$\beta_2 \ge \beta_1$,
\[
   R_{T,\beta_2}(f) - R_{T,\beta_1}(f)
   = \frac{\beta_2-\beta_1}{2\pi}\int_{|t|>T}|F|^2 \;\ge\; 0 .
\]
On the finite compression this reads $A(\beta_2) - A(\beta_1) =
(\beta_2-\beta_1)(I - K_T)$, where $K_T$ is the compression of the Fourier
energy-concentration operator and $0 \le K_T \le I$.
\end{proposition}

Proposition~\ref{prop:beta} is recorded because it rules out a tempting move:
\emph{lowering} $\beta$ cannot raise the exact comparison form, its infimum, or
the least eigenvalue of its compression to a fixed subspace. Inspecting only the
factor $\PsiL - \beta$ inside the integral misses the $\beta\|f\|_2^2$ term. The
identity was checked numerically with maximum residual $5.23\times10^{-143}$ and
diagonal entries of $K_T$ in $[0.9420888637695,\, 0.9980132600038]$ at $T=200$,
$N=16$.

\subsection{Step 3: the Legendre basis}

Let $\varphi_n(u) = \sqrt{(2n+1)/(2L)}\, P_n(u/L)$ be the orthonormal Legendre
basis of $L^2(-L,L)$. Then
\[
   F_n(t) \;=\; i^n \sqrt{2L(2n+1)}\; j_n(Lt) ,
   \qquad \sum_{n \ge 0}(2n+1) j_n(x)^2 = 1 ,
\]
with $j_n$ the spherical Bessel function. The form is assembled in the real
coordinates $y_n = \sigma_n x_n$ with $\sigma_n = (-1)^{\lfloor n/2 \rfloor}$,
which makes $v_n(t) = \sqrt{2L(2n+1)}\,j_n(Lt)$ the real working transform. This
basis is used, rather than the cosine basis of the repository's original
assembly, because $j_n(x)$ is super-exponentially small for $n > ex/2$, which is
what makes Step 4 effective; a cosine basis gives only $O(n^{-2})$ tail decay
and cannot reach the required magnitude at an acceptable dimension.

\subsection{Step 4: the complement}

For $f \in V_N^{\perp}$ with $\|f\|_2 = 1$, Cauchy--Schwarz in $n$ gives
$|F(t)|^2 \le 2L\, S_N(tL)$ with $S_N(x) = \sum_{n \ge N}(2n+1)j_n(x)^2$. Since
$S_N$ is increasing on $0 \le x < N$,
\[
   g_N \;:=\; \sup_{f \in V_N^{\perp},\, \|f\|=1}
     \frac{1}{2\pi}\int_{-T}^{T}|F|^2
   \;\le\; \frac{2LT}{\pi}\, S_N(TL) .
\]
This is a lossy but rigorous surrogate for the $(N+1)$-st Slepian
eigenvalue~\cite{Slepian},
and it is super-exponentially small past the Shannon number $2TL/\pi$. At
$T = 400$, $N = 440$:
\[
   S_N(320) = 9.260051461963608\times10^{-63},
   \qquad
   g_N = 1.886442193224756\times10^{-60}.
\]

\subsection{Step 5: the Schur-type block bound, with strict tests}

Split $f = f_N + f_{\perp}$. With $a$ a certified lower bound for the least
eigenvalue of the compression of $R_{T,\beta}$ to $V_N$, $d$ a certified lower
bound for $R_{T,\beta}$ on $V_N^{\perp}$, and $\kappa$ a certified bound on the
coupling, a $2\times2$ comparison yields positivity as soon as $ad > \kappa^2$.
The threshold $\kappa^2/d$ is a \emph{break-even point}, not the resulting
bound. To avoid the cancellation-prone explicit eigenvalue formula, the final
implication is tested as three \emph{strict} interval signs for the proposed
constant $\mu$:
\begin{equation}
\label{eq:schur}
   a > \mu , \qquad d > \mu , \qquad (a - \mu)(d - \mu) > \kappa^2 .
\end{equation}
Equality anywhere in~\eqref{eq:schur} would give only a nonnegative comparison,
so strictness is required and is checked as a sign of an interval, never of a
midpoint. The coupling constant is
$\kappa = 2|p_N||p_{\perp}| + M\sqrt{g_N}$ with
$M = \sup_{|t| \le T}|\PsiL - \beta|$ and $p$ the pole vector
$p_n = \sqrt{2L(2n+1)}\, i_n(L/2)$ taken in the same $y$-coordinates; $|p|^2 =
\int_{-L}^{L} e^{-u}\,du = 2\sinh L$ exactly. $M$ is obtained from the
monotonicity of $H$ alone, with no covering argument.

\subsection{Verified positive definiteness}
\label{sec:pd}

A naive interval $LDL^{\mathsf T}$ factorization is unusable at this dimension:
entry radii of order $10^{-38}$ amplify by roughly $6.5$ per pivot and reach
$O(1)$ by step $47$. Instead, positive definiteness of $B = A - aI$ is
established by a preconditioned congruence in the sense of Rump~\cite{Rump}. Let $R$ be the
upper triangular Cholesky factor of the \emph{midpoint} of $B$, treated as exact
input, and set $C = R^{-1}$ and
\[
   G \;=\; C^{*} B\, C \;=\; R^{-\mathsf T}\,(A - aI)\, R^{-1},
\]
computed by two explicit forward substitutions in ball arithmetic. The verified
inequality is
\[
   \| I - G \|_{\infty} \;<\; 1 .
\]
This proves $B \succ 0$: the bound places every eigenvalue of the Hermitian $G$
in $(0,2)$, so $G \succ 0$; if $C$ were singular then $G = C^{*}BC$ would be
singular, so $C$ is nonsingular; and then for $x \ne 0$, writing $x = Cy$ with
$y \ne 0$, $x^{*}Bx = y^{*}Gy > 0$. It is a congruence, not a one-sided inverse
residual --- the latter would establish nonsingularity only.

\subsection{Quadrature}

The interior integral in~\eqref{eq:R} is evaluated by panelled
Gauss--Legendre quadrature with certified nodes and weights (FLINT's
\src{legendre\_p\_root}). The symbol, prime comb and the full basis-transform
vector are evaluated once per node and reused across all matrix entries, so no
entry is integrated adaptively. The quadrature remainder is bounded by a
Bernstein-ellipse estimate on genuinely holomorphic extensions: the half-sum
\[
   H_{\mathrm{hol}}(z) \;=\; \tfrac12\left[
     \psi\!\left(\tfrac14 + \tfrac{iz}{2}\right)
   + \psi\!\left(\tfrac14 - \tfrac{iz}{2}\right)\right] - \log\pi
\]
agrees with $H$ on the real axis and has its nearest poles at $z = \pm i/2$
(note that $\operatorname{Re}\psi(\frac14 + \frac{iz}{2})$ is \emph{not} a
holomorphic extension), and the basis product is extended as $v_n(z)v_m(-z)$
rather than by conjugation at a complex argument. On the ellipses, $|v_n(z)|$ is
bounded uniformly in $n$ by
\begin{equation}
\label{eq:cs}
   \Big|\int_{-L}^{L}\varphi(u) e^{izu}\,du\Big|
   \;\le\; \Big(\int_{-L}^{L} e^{-2u \operatorname{Im} z}du\Big)^{1/2}
   \;\le\; \sqrt{2L}\, e^{L|\operatorname{Im} z|}
\end{equation}
for any $L^2$-normalized $\varphi$ supported in $[-L,L]$. Inequality
\eqref{eq:cs} is immediate from Cauchy--Schwarz and is \emph{not} claimed as
new; it is recorded because it replaces a direct interval evaluation of
$j_n$ on wide complex balls, which is numerically useless here. The remainder is
folded into the matrix entries as a widening exactly once.

\subsection{The $T = 400$ certificate}

\begin{center}
\begin{tabular}{ll}
\toprule
parameter & value \\
\midrule
window & $L = 4/5$ \\
cutoff & $T = 400$ \\
dimension & $N = 440$ (Legendre degrees $0..439$; $220$ modes per parity) \\
quadrature & $800$ panels of half-length $1/4$, Gauss--Legendre order $64$ \\
Bernstein data & $b_e = 0.30$ (inside the strip $|\operatorname{Im} z| < 1/2$), $\rho = 2.762049935$ \\
remainder factor & $3.551349410\times10^{-57}$ \\
precision & $900$ bits (Bessel), $384$ bits (accumulation), $1400$ bits (verification) \\
\bottomrule
\end{tabular}
\end{center}

$N = 440$ was chosen from the actual inequalities, not from the heuristic
$eLT/2 = 434.9$. The certified constants are
\begin{align*}
  \beta &= 1.211613695058064549471502, &
  M &= 9.525770639507350587208369, \\
  \kappa &= 1.308343771285388259029018\times10^{-29}, &
  \kappa^2/d &= 1.412796364751588\times10^{-58},
\end{align*}
with pole tail $|p_{\perp}| = 7.993372152149564\times10^{-1282}$.

\begin{center}
\begin{tabular}{lccccc}
\toprule
sector & $a$ (certified) & $\|I-G\|_{\infty}$ & $d$ & $\mu$ & margin in \eqref{eq:schur} \\
\midrule
even & $1.43300048828\times10^{-17}$ & $1.1201203\times10^{-12}$ & $1.211613695058065$ & $1.43300048827856\times10^{-17}$ & $1.74472372088\times10^{-29}$ \\
odd  & $1.3857\times10^{-14}$  & $5.6557966\times10^{-18}$ & $1.211613695058065$ & $1.38569999999861\times10^{-14}$  & $1.68414303613\times10^{-26}$ \\
\bottomrule
\end{tabular}
\end{center}

The whole-window constant is the smaller of the two sector constants, so
$\mu = 1.43300048827856\times10^{-17}$ and \emph{the even sector binds}. Assembly took
$3532$ s, the full run $3581$ s, with peak resident memory $393.4$ MiB; the two
Rump factorizations at $N = 440$ took under $50$ s together.

\subsection{Parity and complex tests}
\label{sec:parity-complex}

The form decouples by parity: the interior matrix has $C_{nm} = 0$ when $n+m$ is
odd, because $v_nv_m$ is then odd on a symmetric interval, and the pole term
$2(C^2 - S^2)$ separates into its even and odd parts. For complex $f = u + iv$
the Hermitian pole term~\eqref{eq:pole} gives
$2(|C|^2 - |S|^2) = 2(C_u^2 - S_u^2) + 2(C_v^2 - S_v^2)$, and the multiplier
term splits because $|F|^2 = |F_u|^2 + |F_v|^2$ pointwise, so
$Q(f) = Q(u) + Q(v)$; each real part then splits by parity. Hence the two sector certificates cover the full class, which is what
makes Theorem~\ref{thm:main}(i) a statement about arbitrary complex $f$.

\begin{remark}
A positive-definiteness check cannot serve as an oracle for the pole sign. The
direct convention checks of Section~\ref{sec:upper} are retained for that
purpose, because an earlier version of this computation carried a pole-phase
error that produced a plausibly negative pivot
(Appendix~\ref{app:errors}, item~1).
\end{remark}

\section{The original-form upper bound}
\label{sec:upper}

\subsection{The certified bound, from the closed-form route}
\label{sec:upperclosed}

A variational upper bound needs one admissible test function and an evaluation
of $Q$ on it that is certified end to end. The evaluation used for
Theorem~\ref{thm:main}(ii) is the \emph{closed-form} one: in the even cosine
basis $f_k(u) = \cos(k\pi u/L)$ the prime sum is finite and exact, and the
archimedean term is summed analytically through the digamma, Lerch $\Phi$ and
Hurwitz $\zeta$ functions, each evaluated as an Arb ball. \textbf{This route
performs no quadrature of any kind}, so it carries no quadrature truncation
error to bound --- which is exactly why it, and not the sharper route of
Section~\ref{sec:timedomain}, supplies the certified endpoint.

The procedure has two stages, which are kept apart. An exploratory
floating-point eigensolver \emph{locates} a candidate vector in this basis; it
contributes no certified claim, and a poor candidate costs a weaker bound, never
a wrong one. The candidate is then frozen to exact rationals and its Rayleigh
quotient enclosed in Arb inside an owned precision scope. Because the second
stage is exact, the first needs only enough precision to find a good trial
function, which is why the mode count can be pushed without an expensive
eigensolve.

Raising the mode count lowers the bound, and the sequence is the useful record:

\begin{center}
\begin{tabular}{rll}
\toprule
$N$ & certified $Q(f_N)/\|f_N\|_2^2 \le$ & \\
\midrule
$28$  & $2.1129926117214908\times10^{-17}$ & (the phase-1 witness) \\
$48$  & $1.9204310132\times10^{-17}$ & \\
$64$  & $1.8541624482\times10^{-17}$ & \\
$96$  & $1.7476544722\times10^{-17}$ & \\
$128$ & $1.6972723502\times10^{-17}$ & \\
$160$ & $1.6731255341\times10^{-17}$ & used in Theorem~\ref{thm:main}(ii) \\
\bottomrule
\end{tabular}
\end{center}

The witness $f_{160}$ is a frozen exact rational vector over the $160$ even modes
$k = 0,\dots,159$ with denominator $10^{180}$. Re-certified from those exact
rationals alone --- no eigensolver, no floating point --- at $2920$ bits,
\[
  \frac{Q(f_{160})}{\|f_{160}\|_2^2}
  \;=\; 1.67312553407068186807726837736\times10^{-17} \pm 4.41\times10^{-47},
\]
whose outward upper endpoint is
\[
  U \;=\; 1.67312553407068186807726837736\times10^{-17}
  \;<\; 1.674\times10^{-17},
\]
giving Theorem~\ref{thm:main}(ii). The radius is a consequence of the route:
with no quadrature and no truncation, only ball rounding contributes. For
comparison, \cite{Zhu} reports a numerical reference estimate
$1.656\times10^{-17}$ for this infimum (Section~\ref{sec:compared}), so the
certified endpoint above sits about $1\%$ over it.

\begin{remark}[on quoting an endpoint]
\label{rem:endpoint}
A serialized enclosure must be consumed through the correct endpoint --- lower
for a lower bound, upper for an upper bound --- and never through its midpoint
or its printed decimal. Two earlier drafts of this note violated that: one
quoted an enclosure's midpoint as the bracket's upper endpoint, and one quoted
an endpoint obtained by re-parsing a stored ball at the ambient $53$ bits. Both
are corrected (Appendix~\ref{app:errors}, item~6), and the safe decimals
$2.1130\times10^{-17}$ here and $1.4330\times10^{-17}$ in part~(i) are the forms
in which the result should be quoted.
\end{remark}

\subsection{A sharper candidate: the time-domain route}
\label{sec:timedomain}

A variational upper bound requires $Q$ itself, not $R_{T,\beta}$. Evaluating
\eqref{eq:Q} directly would mean certifying an oscillatory integral over all of
$\mathbb{R}$. Instead $Q$ is evaluated on the geometric side via the
autocorrelation
\[
   g(x) \;=\; \int f(s+x) f(s)\, ds ,
\]
supported in $[-2L,2L]$, together with the compact archimedean identity
(\cite[Lemma 2.5]{Zhu}, transcribed into the present normalization)
\begin{equation}
\label{eq:arch}
   \mathrm{Arch}(g) = -\big[\gamma_E + \log\pi + \log(1 - e^{-4L})\big] g(0)
   + \int_{0}^{2L} \frac{2\big[e^{-2x}g(0) - e^{-x/2}g(x)\big]}{1 - e^{-2x}}\, dx ,
\end{equation}
\begin{equation}
\label{eq:Qcompact}
   Q(f) = 2sF(i/2)^2 + \mathrm{Arch}(g)
   - \sum_{\log n < 2L} \frac{2\Lambda(n)}{\sqrt n}\, g(\log n) .
\end{equation}
This is the route by which~\cite[\S6.2]{Zhu} evaluates its own variational upper
bounds at this window, there with shift-correlation integrands of degree at most
$239$ integrated exactly by a $160$-node Gauss rule. Two features make it cheap
and rigorous. They are set out below because the present error accounting
depends on them, not as new observations.

\paragraph{Exact inner quadrature.} If $f$ is a polynomial of degree $d$ on
$[-L,L]$ then $g$ restricted to $[0,2L]$ is a polynomial of degree at most
$2d+1$ --- not $2d$: for $f \equiv 1$ one has $g(x) = 2L - x$. What governs the
\emph{inner} rule is the degree of the integrand $f(s+x)f(s)$ in the
integration variable $s$, which is at most $2d$, so an $m$-node
Gauss--Legendre rule is exact once $2m - 1 \ge 2d$, i.e.\ $m \ge d+1$. The implementation uses $m = d+2$; for $d = 438$ that is $440$
nodes. \emph{Exactness} here means no polynomial quadrature truncation error in
exact arithmetic --- node, weight, coefficient and rounding enclosures are still
carried as ball radii, and the outer integral in~\eqref{eq:arch} keeps its own
certified Bernstein remainder.

\paragraph{The removable endpoint, by an identity.} The integrand
of~\eqref{eq:arch} is removable at $x = 0$ and its limit involves
$c = -g'(0^{+})$. For even $f$,
\[
   g'(x) \;=\; -f(L) f(L-x) + \int_{-L}^{L-x} f'(s+x) f(s)\, ds ,
\]
and at $x = 0$ the integral is $\tfrac12[f^2]_{-L}^{L} = 0$, so
\begin{equation}
\label{eq:c}
   c \;=\; -g'(0^{+}) \;=\; f(L)^2 \qquad\text{exactly,}
\end{equation}
and the integrand limit is $c - \tfrac32 g(0)$. This closed form is an
implementation convenience obtained by elementary differentiation and
integration; \cite[Lemma 2.5]{Zhu} leaves $c$ as the one-sided limit and remarks
only that it is nonzero in general. Its absence from the displayed formulas
there is not evidence that it is new, and no novelty is claimed for it. The outer integrand is analytic
in $|\operatorname{Im} x| < \pi$; panels are geometric near the origin so that no
Bernstein ellipse encloses it, and the residual head $[0,10^{-40}]$ is bounded
by an explicit constant.

\subsection{Validation of the representation}

\begin{enumerate}
\item[(a)] On the constant $f = \varphi_0$, the compact route gives
  $Q(\varphi_0) = 0.080274054959269611214 \pm 3.09\times10^{-22}$, matching the
  repository's independent closed form $Q[0,0] = 0.080274054959269611214$ to all
  $21$ printed digits (difference $-3.08\times10^{-22}$). That closed form is the
  evaluator of~\cite{Sensitivity}, which sums the archimedean series analytically
  through the digamma, Lerch $\Phi$ and Hurwitz $\zeta$ functions; it is an
  independent route to $Q$, not an independent implementation of the present
  one.
\item[(b)] A degree-$0$ check cannot see a pole-phase error, since
  $\sigma_0 = +1$. Checks at $\varphi_2$ and at the signed mixed function
  $\varphi_0 - \varphi_2$ were therefore added, with the pole moment computed
  \emph{directly} from the physical-space function rather than from the
  pole-vector helper, and with the finite prime shifts re-enumerated
  independently of the symbol helper:
  \begin{center}
  \begin{tabular}{lll}
  \toprule
  $f$ & $Q(f)$ & pole moment (direct) \\
  \midrule
  $\varphi_0$ & $0.080274054959269611214$ & $1.29891290374844725$ \\
  $\varphi_2$ & $0.389336213205013655883$ & $0.030516224225381024$ \\
  $\varphi_0 - \varphi_2$ & $0.117009437226589934780$ & $1.26839667952306623$ \\
  \bottomrule
  \end{tabular}
  \end{center}
  All agree with the helper values, the first to $10^{-25}$.
\item[(c)] \emph{Disclosure.} This is a second \emph{representation} of the same
  form, not an independent implementation: it shares the Arb library, the symbol
  helper for the comb data, and the pole-vector helper. It is not described as an
  independent check.
\end{enumerate}

\subsection{The frozen test functions of the time-domain route}

Two trial functions are used, each a frozen exact rational vector in the
Legendre $y$-coordinates with denominator $10^{40}$:
\begin{itemize}
\item $f_{200}$: $110$ even modes, degrees $0..218$, $\|f\|_2^2 =
  2.275244443650965360970550$;
\item $f_{400}$: $220$ even modes, degrees $0..438$.
\end{itemize}
Both were \emph{located} by an exploratory floating-point eigensolver, then
rounded to exact rationals and certified; the eigensolver contributes no
certified claim, and in particular supplies no eigenvalue.

For $f_{200}$, equation~\eqref{eq:Qcompact} gives, at $94$ outer panels and head
cut $10^{-40}$,
\[
  \text{pole} = 3.2303213242823875400,\qquad
  \mathrm{Arch} = -3.1402989129419417722,
\]
\[
  \text{comb} = 0.090022411340445684793,
\]
whose signed sum is $8.30143811294573721\times10^{-17}$: a cancellation of about
sixteen digits, which is why the evaluation is done at high precision in ball
arithmetic and why the compact representation --- with its exact inner
quadrature and the exact identity~\eqref{eq:c} --- matters. Dividing by the norm,
\[
  q(f_{200}) = 3.64859175290407702\times10^{-17} \pm 4.15\times10^{-35}.
\]
This is above the then-current upper bound and so does not improve it; the
expected outcome, since $f_{200}$ minimizes a comparison form, not $Q$.

For $f_{400}$,
\[
  q(f_{400}) = 1.90138291114546440\times10^{-17} \pm 6.03\times10^{-35},
\]
which would improve on $U$ --- were the enclosure complete. It is not; see
Section~\ref{sec:status}. The radii printed above reflect ball rounding and the
explicit head bound only.

\subsection{Excluded route: the time-domain evaluator}
\label{sec:status}

The outer integral in~\eqref{eq:arch} is evaluated by panelled Gauss--Legendre
quadrature --- geometric panels from $\eta = 10^{-40}$ to $10^{-3}$, then $32$
uniform panels of order $32$ to $2L$ --- and the implementation adds the
explicit head bound on $[0,\eta]$ but \textbf{no bound on the quadrature
truncation error of those panels}. Ball arithmetic carries the node, weight and
rounding enclosures; it does not capture the difference between a
Gauss--Legendre sum and the integral it approximates, which is a property of the
rule rather than of the arithmetic. The radii printed in
Section~\ref{sec:timedomain} therefore do not include that term.

Accordingly, \textbf{every value produced by the time-domain route is withdrawn
as a certificate} pending such a bound: $q(f_{200})$, $q(f_{400})$, the endpoint
once quoted from the latter, and the discarded energies
$e_T = q - r_T$ of Section~\ref{sec:cutoff} together with their ratios. They are
retained below as \emph{measurements}, labelled as such.

What is \emph{not} affected, and why:
\begin{itemize}
\item the lower bound of Section~\ref{sec:lower}, whose interior quadrature
  carries a certified Bernstein-ellipse remainder folded into the matrix entries
  exactly once, re-derived by its verifier;
\item the odd-sector bound of Theorem~\ref{thm:main}(v), from the same matrix route;
\item the certified upper bound of Section~\ref{sec:upperclosed}, which uses no
  quadrature at all;
\item the scalar values $R_T(f)$ of Section~\ref{sec:cutoff}, which are computed
  by a separate routine that \emph{does} add a certified remainder and publishes
  it; in particular the ceiling $\lR(800) \le r_{800}(f_{400})$ does not involve
  $Q$ at any point.
\end{itemize}

A rigorous remainder for~\eqref{eq:arch} requires a certified bound on the
analytic continuation of the outer integrand on a Bernstein ellipse around each
panel, with the denominator's zeros at $x \in i\pi\mathbb{Z}$ kept outside. The
degree statement corrected above is a prerequisite for that bound and not a
substitute for it: knowing $\deg g \le 2d+1$ says nothing by itself about the
size of $g$ off the real axis. The obvious route --- bound $|g|$ on the ellipse
by $\max_{[0,2L]}|g| \cdot R^{\deg g}$ --- is far too lossy at $\deg g = 877$,
forcing panels of width $\sim10^{-3}$. A coefficient-sensitive bound should
work instead, because the frozen minimizers concentrate on \emph{low} Legendre
orders, so estimating $|f(w)|$ off the real axis directly from the frozen
coefficients is tight exactly where $R^{\deg g}$ is hopeless. That derivation is
left open; it is recorded with the diagnostic in
\src{phase5/src/outer\_convergence.py}, which measures the refinement difference
at $\pm 2.71\times10^{-35}$ --- small, and not a bound.

Since the certified endpoint of Section~\ref{sec:upperclosed} now sits
\emph{below} the withdrawn time-domain value, closing this gap would no longer
improve the bracket; it would restore a second, independent route to it.

\section{The parity separation}
\label{sec:parity}

The $T = 400$ certificate of Section~\ref{sec:lower} delivers
$a = 1.3857\times10^{-14}$ in the odd sector against
$a = 1.43300048828\times10^{-17}$ in the even sector.
Re-evaluating the three strict signs~\eqref{eq:schur} from the exact certified
data --- not from rounded report values --- at $\mu_{\mathrm{odd}} =
1.3856\times10^{-14}$ gives
\[
   a > \mu_{\mathrm{odd}},\qquad d > \mu_{\mathrm{odd}},\qquad
   (a - \mu_{\mathrm{odd}})(d - \mu_{\mathrm{odd}}) > \kappa^2 ,
\]
with margin $1.68414303613\times10^{-26}$. So the whole odd sector satisfies
$Q(f) \ge 1.38569999999861\times10^{-14}\|f\|_2^2$. Combining with the certified endpoint $U_0$
of Section~\ref{sec:upperclosed}, and using that the witness $f_{160}$ is even,
\[
   \mu_{\mathrm{odd}} - U_0 \;=\; 1.384026874464\times10^{-14}
   \;>\; 0 ,
\]
so
\[
   \inf_{\text{even}} \frac{Q(f)}{\|f\|_2^2} \;\le\; U_0 \;<\;
   \mu_{\mathrm{odd}} \;\le\; \inf_{\text{odd}} \frac{Q(f)}{\|f\|_2^2},
\]
which with the orthogonal parity decomposition identifies the even infimum as
$\lQ(4/5)$. The global lower bound is $1.43300048827856\times10^{-17}$, governed by the
even sector.

The witness used here is the closed-form one. An earlier draft ran this argument
on the time-domain value for $f_{400}$; that value is withdrawn by
Section~\ref{sec:status}, and the separation is re-established above on a
witness whose evaluation uses no quadrature. The conclusion is unchanged.

As noted in Section~\ref{sec:prior}, none of this \emph{structure} is new: at
this very window \cite[Thm.~6.2]{Zhu} certifies the same ordering, with
$\lambda^{\mathrm{odd}}_1 \ge 8.206\times10^{-15}$. The odd constant recorded
above exceeds that value by a factor $1.68$, which is a sharpening of one
constant and nothing more; the separation itself, and the simplicity and
evenness of the ground state at this window, were established
by~\cite[Thm.~6.2]{Zhu} and are not reproduced here. Nothing about simplicity,
uniqueness or attainment follows from the numbers above.

\section{Cutoff diagnostics}
\label{sec:cutoff}

\subsection{The cutoff comparison, proved}

\begin{proposition}
\label{prop:cutoff}
Let $\PsiL \ge \beta_i$ on $|t| > T_i$ for $i = 1,2$, with $T_2 \ge T_1$ and
$\beta_2 \ge \beta_1$. Then
\[
   R_{T_2,\beta_2}(f) - R_{T_1,\beta_1}(f)
   = \frac{1}{2\pi}\int_{T_1 < |t| \le T_2} (\PsiL - \beta_1)|F|^2
   + \frac{\beta_2 - \beta_1}{2\pi}\int_{|t| > T_2} |F|^2
   \;\ge\; 0 ,
\]
both terms being nonnegative. Hence $\lR(T_2) \ge \lR(T_1)$ for the exact forms.
\end{proposition}

Two cautions. This does not imply that a \emph{certified numerical} bound
improves, since error terms change with the parameters. And it does not order
finite matrices of different sizes: changing $N$ changes the compression, so
only matrices in the same subspace are comparable by this identity.

\subsection{The scalar table}

Each $R_T$ entry below is a certified enclosure on a \emph{fixed, frozen}
function; no reoptimization is involved, so the columns compare exact forms on
the same vector. The $Q$ row comes from the time-domain route and is therefore a
\textbf{measurement, not a certificate} (Section~\ref{sec:status}); it is shown
because the diagnostics below are read off it.

\begin{center}
\begin{tabular}{lll}
\toprule
 & $f_{200}$ & $f_{400}$ \\
\midrule
$R_{200}$ & $1.0310220302459896\times10^{-17}$ & $1.24275215\times10^{-17}$ \\
$R_{400}$ & $2.3520530467310079\times10^{-17}$ & $1.4330724047355850\times10^{-17}$ \\
$R_{800}$ & $3.01603\times10^{-17}$            & $1.6681\times10^{-17}$ \\
$Q$ \tier{measured} & $3.6485917529040770\times10^{-17}$ & $1.9013829111454644\times10^{-17}$ \\
\bottomrule
\end{tabular}
\end{center}

The $R_{800}$ entries carry wider enclosures ($\pm 5.30\times10^{-23}$ and
$\pm 4.60\times10^{-22}$) because the quadrature remainder grows with the
cutoff; $R_{200}(f_{400})$ was evaluated as a direct scalar integral on the full
degree-$438$ function, with no truncation and no enlarged matrix. Each column is
strictly increasing in $T$, as Proposition~\ref{prop:cutoff} requires.

\subsection{What the comparison discards}

From~\eqref{eq:e}, with $e_T(f) = q(f) - r_T(f)$ by rigorous interval
subtraction. Since $q$ is withdrawn as a certificate
(Section~\ref{sec:status}), \textbf{every entry in this table is a measurement},
and no statement below rests on one:
\begin{center}
\begin{tabular}{lll}
\toprule
 & $f_{200}$ & $f_{400}$ \\
\midrule
$e_{200}$ & $2.617569722658087\times10^{-17}$ & $6.586308\times10^{-18}$ \\
$e_{400}$ & $1.296538706173069\times10^{-17}$ & $4.683105064098794\times10^{-18}$ \\
$e_{800}$ & $6.3256\times10^{-18}$            & $2.333\times10^{-18}$ \\
\midrule
$e_{400}/e_{200}$ & $0.49532156 \pm 4.69\times10^{-9}$ & $0.7110365 \pm 1.61\times10^{-8}$ \\
$e_{800}/e_{400}$ & $0.48788 \pm 9.07\times10^{-6}$    & $0.498 \pm 2.36\times10^{-4}$ \\
\bottomrule
\end{tabular}
\end{center}

\textbf{There is no decay law here.} On $f_{400}$ the successive ratios are
$0.711$ and then $0.498$, which is not a constant rate; and the two functions
disagree at the \emph{same} doubling $200 \to 400$, giving $0.495$ and $0.711$,
so the ratio is not a function of $T$ alone. That both functions give
approximately $0.49$ across $400 \to 800$ is two samples of one step, not a rate
in $T$, and comparing differently optimized vectors would establish less still.
No uniform operator estimate is inferred from a fixed-vector ratio, and none of
these ratios is used in any certified statement.

\subsection{What $q_f - r_f$ is, and is not}

For $f_{400}$, $q_f - r_f = 4.6831050640988\times10^{-18}$ is the
discarded-term identity~\eqref{eq:e} evaluated on that function, inside its
enclosure --- a measured quantity, since $q$ is withdrawn
(Section~\ref{sec:status}). It is \emph{not} $\lQ - \lR(400)$: both $\lR(400) \le r_f$
and $\lQ \le q_f$ are upper bounds, and two upper bounds do not make their
difference an equality between infima.

\subsection{Certified ceilings, and the stopping decision}

Because $f_{400}$ and $f_{200}$ are admissible, each $R_T$ value in the table is
a variational \emph{upper} bound for $\lR(T)$, and any certificate resting only
on $Q \ge R_{T,\beta}$ is capped by it. Relative to the lower bound in force:
\begin{center}
\begin{tabular}{llll}
\toprule
route & certified ceiling & reference lower bound & headroom \\
\midrule
$R_{200}$ & $\lR(200) \le r_{200}(f_{200}) \le 1.0310221\times10^{-17}$ & $1\times10^{-17}$ & $3.1023\%$ \\
$R_{400}$ & $\lR(400) \le r_{400}(f_{400}) \le 1.4330725\times10^{-17}$ & $1.4330005\times10^{-17}$ & $0.0050\%$ \\
$R_{800}$ & $\lR(800) \le r_{800}(f_{400}) \le 1.6681461\times10^{-17}$ & $1.4330005\times10^{-17}$ & $16.409\%$ \\
\bottomrule
\end{tabular}
\end{center}

Each ceiling above is the \emph{outward-rounded upper endpoint} of the
corresponding enclosure, not its midpoint: from
$r_{800}(f_{400}) = 1.6681\times10^{-17} \pm 4.60\times10^{-22}$ the endpoint is
$1.6681460\times10^{-17}$, quoted outward as $1.6681461\times10^{-17}$; quoting
the midpoint $1.6681\times10^{-17}$ as a ceiling would be the same error as in
Remark~\ref{rem:endpoint}.

So the $T = 400$ route is within about $1.64\%$ of exhausted, and the entire
$T = 800$ route is capped at a factor $1.1831$ over the present lower endpoint;
since a reoptimized $f_{800}$ can only sit at or below that ceiling, the
realizable gain is at most that. Measured costs on the machine used:

\begin{center}
\begin{tabular}{ll}
\toprule
stage & measurement \\
\midrule
assembly, $T=200$, $N=220$, $800\times64$ nodes & $1338$ s, peak $205.5$ MiB \\
assembly, $T=400$, $N=440$, same node count    & $3532$ s, peak $393.4$ MiB \\
Rump factorizations, both sectors, $N=440$     & under $50$ s \\
compact $Q(f_{200})$                           & $389$ s \\
cross table ($R_{400}$, $Q(f_{400})$)          & $1544$ s \\
scalar leg $R_{200}(f_{400})$, $12800$ nodes   & $933$ s \\
scalar leg $R_{800}(f_{400})$, $51200$ nodes   & $8049$ s \\
scalar leg $R_{800}(f_{200})$, $51200$ nodes   & $3651$ s \\
\bottomrule
\end{tabular}
\end{center}

The assembly grew by a factor $2.64$ when $N$ doubled at fixed node count, and
the $T = 800$ scalar legs overran the node-count scaling by more than a factor
of two, which indicates that per-Bessel cost grows with the argument rather than
node count alone. Extrapolating --- as an estimate anchored on these
measurements, not a guarantee --- a full $T = 800$ matrix at $N \approx 880$ with
roughly $1600\times64$ nodes would cost on the order of ten hours of assembly.
Against a ceiling of $+18.31\%$, the recommendation recorded in the
accompanying repository is to \textbf{consolidate at $T = 400$}. Neither the
pilot nor any failed certificate here is an intrinsic barrier; this is a
cost--benefit judgement from measured data, and a separately scoped $T = 800$
attempt would be a reasonable thing to run with that ceiling stated in advance.


\section{The second window}
\label{sec:second}

Everything in Sections~\ref{sec:lower} and~\ref{sec:upperclosed} is run a second
time at
\[
   \widetilde L \;=\; \frac{80471895621705018}{10^{17}}
   \;=\; 0.80471895621705018 ,
\]
a rational chosen strictly below $L_\ast = (\log 5)/2 = 0.8047189562170501873\ldots$,
with $L_\ast - \widetilde L < 7.3\times10^{-18}$. The choice matters for cost,
not for the statement. The prime comb of~\eqref{eq:Psi} runs over $n$ with
$\log n < 2L$, and $2\widetilde L < \log 5$, so $n = 5$ is excluded and
\[
   A_{\widetilde L} \;=\; A_{4/5} \;=\; 2.941973525223620455504 ,
\]
whence the exterior gate of Lemma~\ref{lem:gate} admits the same cutoff
$T = 400$ (the gate threshold is $T_{\mathrm{gate}} = 119.0869\ldots$ at both
windows). Had $\widetilde L$ been taken above $L_\ast$, the comb would acquire
$n = 5$, $A_L$ would jump by $2\log 5/\sqrt 5$ and $T_{\mathrm{gate}}$ would rise
to $502.39\ldots$, multiplying the assembly cost. At $L_\ast$ itself the strict
inequality still excludes $n = 5$; the jump happens for $L > L_\ast$.

With $N = 440$, $800$ panels of half-length $1/4$ and Gauss--Legendre order
$64$, the certified constants are
\[
   \beta_{400} = 1.211613695058064549472, \quad
   \kappa = 7.756529961646111179258\times10^{-29}, \quad
   d = 1.211613695058064549472 ,
\]
and the two sector shifts, each verified by the preconditioned congruence of
Section~\ref{sec:pd} on the assembled block, are
\[
   a_{\mathrm{even}} = 8.27343868255615234375\times10^{-18}, \qquad
   a_{\mathrm{odd}} = 7.885009765625\times10^{-15} ,
\]
with $\|I-G\|_\infty$ equal to $5.0567\times10^{-11}$ and
$4.0104\times10^{-16}$ respectively. The strict Schur test then yields
Theorem~\ref{thm:main}(iii). The upper bound (iv) is the closed-form Rayleigh
quotient of Section~\ref{sec:upperclosed} at a frozen rational vector in the
same $160$-mode even cosine basis, at the new window.

\begin{remark}
$\widetilde\lambda^2 = e^{2\widetilde L}$ agrees with $5$ to about
$7.3\times10^{-17}$; the decimal $4.999999999999999927$ is a rounding of that
quantity, not its definition. The definition is the displayed rational
$\widetilde L$.
\end{remark}

\section{Conditioned coercivity, by parity sector}
\label{sec:cond}

The smallness of $\lQ(4/5)$ is carried by a few directions. This is made
precise, and certified, by a rank-$K$ device that needs no new assembly.

\begin{lemma}
\label{lem:bump}
Let $A$ be the compression of $R_T$ to $V_N$ in an orthonormal basis, let
$v_1,\dots,v_K$ be any vectors and $c \ge 0$. If
$A + c\sum_{i} v_i v_i^{\!\top} \succeq a I$ then
$\langle A f, f\rangle \ge a\|f\|^2$ for every $f$ with
$\langle v_i, f\rangle = 0$, $i = 1,\dots,K$.
\end{lemma}

\begin{proof}
On that subspace the rank-$K$ term contributes exactly $0$.
\end{proof}

Because the implication holds for \emph{any} $v_i$, their provenance needs no
rigour: only the interval test does. The vectors used here were located by
inverse iteration on the midpoint matrix and then frozen as exact rationals with
denominator $10^{34}$; they are published in the certificates. Combining
Lemma~\ref{lem:bump} with the unchanged complement and coupling constants
$d$ and $\kappa$ of Section~\ref{sec:lower} through the same strict Schur test
gives, at $L = 4/5$:

\begin{theorem}
\label{thm:cond}
There are explicit rational vectors in the $220$-mode even and odd Legendre
coordinate bases of the accompanying package such that, for admissible $f$
supported in $[-4/5,4/5]$ and orthogonal to the listed directions,
\[
\begin{array}{llr}
\text{even sector,} & 4 \text{ directions:} & Q(f) \ge 0.617553710936882\,\|f\|_2^2 \\
\text{even sector,} & 5 \text{ directions:} & Q(f) \ge 0.853492736815556\,\|f\|_2^2 \\
\text{odd sector,}  & 3 \text{ directions:} & Q(f) \ge 0.106805801391493\,\|f\|_2^2 \\
\text{odd sector,}  & 4 \text{ directions:} & Q(f) \ge 0.784130096434765\,\|f\|_2^2
\end{array}
\]
Over both sectors together, $4$ even $+\,3$ odd conditions give
$Q(f) \ge 0.106805801391493\,\|f\|_2^2$, and $4$ even $+\,4$ odd give
$Q(f) \ge 0.617553710936882\,\|f\|_2^2$.
\end{theorem}

Short forms must round \emph{down}: $0.6175$ is a valid abbreviation of the
four-condition even constant and $0.6176$ is not.

\begin{remark}
These are per-sector statements, and four conditions on the combined space
cannot give the even constant: by min--max the odd sector carries minuscule
eigenvalues of its own, so the combined-space bound with $7$ conditions is the
smaller of the two sector constants.
\end{remark}

\begin{remark}
\label{rem:notprolate}
Removing directions changes the minimum over what remains; it does \emph{not}
raise $\lQ(4/5)$, which stays inside the bracket of Theorem~\ref{thm:main}. And
the conditioning spaces here are the spans of explicitly listed numerical
vectors. They are \emph{not} identified with the prolate conditioning spaces
of~\cite[\S5.3]{ACLG}; that identification would need a separate proof or a
separate certificate, and neither is offered here.
\end{remark}

\section{An eigenvalue count at $L = 4/5$}
\label{sec:count}

Theorem~\ref{thm:cond} bounds a count from above. By Courant--Fischer, four
conditions leaving the infimum at least $0.6175$ means at most four even
eigenvalues of $Q$ lie below $0.6175$, in particular at most four below
$10^{-2}$; the three-condition odd constant gives at most three odd. For the
other direction one needs a test \emph{space}, and the closed-form evaluator
supplies one.

\begin{proposition}
\label{prop:count}
Let $M_N$ be the matrix of $Q$ on the first $N$ even cosine modes in an
orthonormal basis, computed in closed form as in Section~\ref{sec:upperclosed}.
Since $M_N$ is a compression of $Q$, min--max gives
$\lambda_j(M_N) \ge \varepsilon_{\mathrm{even}}(j)$ for every $j$. Hence for any
$K$-dimensional $W \subset \mathbb{R}^N$,
\[
   \varepsilon_{\mathrm{even}}(K) \;\le\; \lambda_K(M_N)
   \;\le\; \lambda_{\max}\!\big(M_N|_W\big) ,
\]
and $\lambda_{\max}(M_N|_W) \le c$ is certified by the $K \times K$ test
$c\,G - B \succ 0$ with $G = W^{\!\top}W$ and $B = W^{\!\top}M_N W$.
\end{proposition}

At $N = 160$, working precision $2920$ bits, with $W$ a frozen rational
$4$-dimensional space (denominator $10^{30}$; the Gram matrix $G$ is verified
positive definite, so $\dim W = 4$), the test certifies
\[
   \varepsilon_{\mathrm{even}}(4) \;\le\; 0.00311839891820
   \;<\; 0.003119 \;<\; 10^{-2} .
\]
Together with the four-condition bound:

\begin{theorem}
\label{thm:exactly4}
At $L = 4/5$, exactly four eigenvalues of $Q$ in the even sector lie below
$10^{-2}$.
\end{theorem}

In the odd sector only \textbf{at most three} is certified: the closed-form
evaluator implements the even cosine basis only, so no odd compression of $Q$ is
available here and the matching lower side is absent. Counts observed in the
compression of the comparison form $R_T$ are \emph{numerical}: since
$Q \ge R_T$, the eigenvalues of a compression of $R_T$ bound those of $Q$ from
neither side for $j \ge 2$. Neither count is transferred to $\widetilde L$ or to
$L_\ast$; that would need its own certificate.

\section{Normalization dictionary, and relation to Andrews--Connes--La~Gatta}
\label{sec:dict}

\cite{ACLG} studies the Weil quadratic form on the multiplicative window
$[\lambda^{-1},\lambda] \subset \mathbb{R}^\ast_+$ with $d^\ast u = du/u$, in the
normalization $\Psi = W_{0,2} - W_{\mathbb{R}} - \sum_p W_p$ of their
equations~(6)--(8). Writing $t = \log u$ and $R(t) = f(t)$, that form and the
form~\eqref{eq:Q} of Section~\ref{sec:setup} are the same object:

\begin{proposition}[Dictionary]
\label{prop:dict}
Let $g(u) = f(\log u)$ and $\lambda = e^{L}$. Then
$\|g\|^2_{L^2(d^\ast u)} = \|f\|^2_{L^2(dx)}$ and
$QW_\lambda(g,g) = Q(f)$, with the two pole signs matching sector by sector.
\end{proposition}

\begin{proof}
In the additive variable, $F(e^s) = C(s) = \int R(t)R(t+s)\,dt$. Their
$W_{0,2}(F) = 2\big(\int R e^{u/2}\big)\big(\int R e^{-u/2}\big)$, which is the
pole term~\eqref{eq:pole} with $s = +1$ for even $R$ and $s = -1$ for odd $R$,
since then $\int R e^{u/2} = -\int R e^{-u/2}$. Their $W_p(F)$ is term by term
the prime sum of~\eqref{eq:Psi}. For the archimedean part, their Lemma 8.13
writes $-W_{\mathbb{R}}$ as the Fourier multiplier $m(\tau) - c_0$ with
$m(\tau) = \int_0^\infty (1-\cos\tau s)\,e^{s/2}/\sinh s\;ds$ and
$c_0 = \log 4\pi + \gamma_E + \int_0^\infty (e^{s/2}-1)/\sinh s\;ds$, so
\[
   m(\tau) - c_0 \;=\; \int_0^\infty \frac{1 - e^{s/2}\cos\tau s}{\sinh s}\,ds
   \;-\; \log 4\pi \;-\; \gamma_E .
\]
From $\psi(z) = -\gamma + \int_0^\infty (e^{-u}-e^{-zu})/(1-e^{-u})\,du$ at
$z = \tfrac14 + \tfrac{i\tau}{2}$, substituting $u = 2s$ and using
$1 - e^{-2s} = 2e^{-s}\sinh s$,
\[
   \Re\,\psi\!\left(\tfrac14 + \tfrac{i\tau}{2}\right)
   \;=\; -\gamma + \int_0^\infty \frac{e^{-s} - e^{s/2}\cos\tau s}{\sinh s}\,ds .
\]
The difference of the two displays is
$\int_0^\infty (1-e^{-s})/\sinh s\;ds - \log 4$, and
$(1-e^{-s})/\sinh s = 2e^{-s}/(1+e^{-s})$ integrates to $2\log 2 = \log 4$.
Hence $m(\tau) - c_0 = \Re\,\psi(\tfrac14+\tfrac{i\tau}{2}) - \log\pi$, which is
the multiplier $\PsiL$ of~\eqref{eq:Psi} up to the prime comb already matched.
\end{proof}

Proposition~\ref{prop:dict} is an \emph{analytic identity}. It was additionally
cross-checked numerically: an independent implementation of~(6)--(8), sharing no
code with this package, reproduces the certified values of $Q$ at the Legendre
test functions $\varphi_0$, $\varphi_2$ and $\varphi_0 - \varphi_2$ to about
thirty digits, matching the stored enclosures to their serialized width. The
cross-check corroborates the identity; it does not replace the proof.

Under the dictionary our window $L = 4/5$ is $\lambda^2 = 4.95303242439511480365$
and $\widetilde L$ is $\lambda^2 = 5 - 7.3\times10^{-17}$, just below the
smallest window tabulated in~\cite{ACLG}.

\paragraph{What is comparable, and what is not.}
The eigenvalues reported in~\cite[\S4]{ACLG} are eigenvalues of finite
log-Fourier compressions; by min--max they bound the limiting eigenvalues from
above only, as those authors state. Their value at $\lambda^2 = 5$, recovered
from a printed $-\log\varepsilon$, is $\approx 9.1663\times10^{-18}$; it lies
inside the bracket of Theorem~\ref{thm:main}(iii)--(iv) at our nearby window.
That placement is \textbf{numerical consistency} between a reconstructed
finite-compression estimate and a certified bracket at a slightly different
parameter --- not a certified comparison at identical parameters, and not an
independently known value of the infinite-dimensional infimum.

Their Hypothesis~3.4 compares $\varepsilon_\lambda(j)$ with the
uncertainty-operator rungs $\delta_{4j}(\lambda)$ up to constants and a factor
$\log\lambda$. It does \emph{not} assert that the ratio of two Weil minima
equals the ratio of the corresponding rungs; that would need the proportionality
factor to be the same at both windows, which is a strictly stronger
approximation. In particular the rung ratio
$\delta_4/\delta_4 = 0.579867\ldots$ between our two windows is a
\textbf{numerical rung-ratio reference}, not a prediction, and it happens to lie
inside the interval of Corollary~\ref{cor:ratio}. All comparisons against
$\delta_{4j}$ in the accompanying report are labelled numerical, because the
rung values there are high-precision rather than interval enclosures.

\section{Reproducibility and limitations}
\label{sec:repro}

\subsection{Artifacts and verification}

Every decisive quantity in this note is published as a midpoint--radius
enclosure in machine-readable form, together with the exact rational input
vectors and the code that produced them. Printed decimals in this note are
summaries; the enclosures are authoritative.

Two layers of checking ship with the package. The \emph{repository suite} ---
eight scripts, $176$ checks --- recomputes the phase-1 to phase-6 chain from the
frozen artifacts and exits nonzero on any failure. On top of it the
\emph{release gate} \src{phase7/acceptance.py} re-verifies every claim of this
note from the packaged files rather than from values held in memory: the
repository suite itself; the closed-form upper bound re-certified from its
stored exact rationals; a Rump replay of \emph{both} improved unconditioned
$L = 4/5$ shifts, which the phase-4 verifier does not cover because it tests the
earlier committed shifts and because the strict Schur test assumes $A \succeq aI$
rather than checking it; all eight Schur checks from their own stored
$a, d, \kappa, \mu$, with rejection required at $\mu = a$; the four conditioned
rank-$K$ bumps rebuilt from the stored rational vectors; both second-window
blocks; the second-window upper quotient; the count certificate with its Gram
matrix; the directed-decimal checks and the ratio of
Corollary~\ref{cor:ratio}; and the provenance and remainder fields the chain
depends on. On the run accompanying this version the gate reports
\textbf{$61$ checks, $0$ failed, in $207$ seconds}, and the repository suite
inside it reports $176$ passing checks.

All of this is \emph{frozen-artifact replay}, not a cold rebuild: the three
matrix assemblies are not re-run (measured at $3532\,\mathrm{s}$ for $T = 400$ at
$L = 4/5$, $3678\,\mathrm{s}$ at $\widetilde L$, $1343\,\mathrm{s}$ for
$T = 200$), and the scripts that regenerate them are named in the package
instructions. The verifiers are not independent implementations: they share the
assembly and constants modules with the computation they check, and no
independence is claimed for them. One script re-derives every decimal quoted
above from the artifacts and additionally checks each quoted bound's
\emph{rounding direction} and its scope and parity label; it is what caught the
mis-quoted upper endpoint of Appendix~\ref{app:errors}, item~6.

Within those scripts two kinds of check must be distinguished. Most recompute a
\emph{certificate-producing inequality} --- a Schur sign, a Cholesky residual, a
tail bound --- in ball arithmetic, and a failure there invalidates a claim.
A few are \emph{auxiliary}: they compare the Arb assembly against the
repository's independent \src{mpmath} implementation, or assert that a
diagnostic spectrum is finite. Those use floating point and settle nothing about
the certificates; one of them formerly compared two matrices after converting
both to ordinary doubles while asserting agreement below $10^{-30}$, a threshold
that no double-precision comparison can reach. It now compares at full precision
and states the tolerance its inputs actually support
(Appendix~\ref{app:errors}, item~9).

The verifiers fail closed. Any NaN, infinite radius, indeterminate enclosure, or
unresolved sign produces an explicit failure or a refinement request; no skipped
cell, panel or pivot is interpreted as success. This matters because an earlier
version of the computation silently skipped cells where a library digamma call
returned NaN on wide balls and thereby produced a wrong constant that was
\emph{used} (Appendix~\ref{app:errors}, item~2).

\subsection{Environment}

Linux, Python $3.14.7$, \src{mpmath 1.3.0}, \src{python-flint 0.9.0} (FLINT/Arb).
Floating-point \src{mpmath} is used to \emph{locate} candidate vectors, which
are then frozen as exact rationals and certified in Arb, and in the auxiliary
cross-implementation checks described above. It enters no
certificate-producing inequality; it does appear in a small number of auxiliary
verifier assertions, which is a weaker statement than the one made in earlier
versions of this note. Arb working precisions on sign-critical paths range from
$256$ to $2000$ bits and are stated per stage in the artifacts.

\subsection{Limitations}

\begin{enumerate}
\item \textbf{Two windows.} Everything is at $L = 4/5$ and $L = \widetilde L$.
  There is no propagation theorem and nothing follows at any window above
  $\widetilde L$ except the upper bound of Proposition~\ref{prop:mono}. In
  particular \emph{no positive lower bound at $L_\ast = (\log 5)/2$} follows
  from these certificates.
\item \textbf{An excluded route, still open.} The time-domain evaluator of
  Section~\ref{sec:timedomain} omits the truncation error of its outer panelled
  rule, so every value it produced remains uncertified
  (Section~\ref{sec:status}) and is excluded from every claim above. A rigorous
  outer remainder is still open; it is off the critical path, since closing it
  would restore a second route to a bracket this note already has.
\item \textbf{No odd count from below.} Theorem~\ref{thm:exactly4} is
  even-sector only. The closed-form evaluator implements the even cosine basis,
  so there is no odd compression of $Q$ here and only ``at most three'' is
  certified in the odd sector.
\item \textbf{The conditioning spaces are not identified with prolate spaces.}
  See Remark~\ref{rem:notprolate}. The constants of Theorem~\ref{thm:cond} hold
  for the explicitly listed rational vectors; no identification
  with~\cite[\S5.1--5.3]{ACLG} is proved or certified.
\item \textbf{Comparisons against the uncertainty-operator rungs are
  numerical.} The rung values $\delta_{4j}(\lambda)$ used in the accompanying
  report are high-precision, validated against the printed values
  of~\cite{ACLG}, but are \emph{not} interval enclosures. Every ratio against
  them is labelled numerical, and none enters a certified claim.
\item \textbf{Rational windows only.} The assembly is parameterized by a
  rational $L$; $L_\ast = (\log 5)/2$ is irrational and cannot be used
  directly, which is why the second window is the rational $\widetilde L$ just
  below it.
\item \textbf{An imported normalization.} The explicit-formula pairing in the
  convention of Section~\ref{sec:setup} is imported from the accompanying
  repository, not rederived there. It is now additionally supported by
  Proposition~\ref{prop:dict}, which identifies it with an independently
  published normalization and is corroborated numerically to about thirty
  digits; that is evidence for the convention, not a derivation of it.
\item \textbf{An unverified external matrix step.} Zhu describes supplementary
  source code and certificates in~\cite[\S17]{Zhu}; these were not obtained or
  executed here, so the matrix step of~\cite[Thm.~1.1]{Zhu} is not reproduced.
  The present certificates are independent of it.
\item \textbf{Shared arithmetic.} The verifiers share library code with the
  computation they check, and are reported as replays, not as independent
  implementations.
\item \textbf{No arithmetic content.} The only analytic sharpening proved here
  is the monotone gate of Lemma~\ref{lem:gate}, which moves a constant by under
  one percent and does not touch the doubly exponential growth of
  $T_{\mathrm{gate}}(L) = 2\pi e^{A_L}(1+o(1))$ that governs the cost in $L$.
  That growth, and not the finite block, is what caps this method.
\item \textbf{Scope.} No zero-free region, no verified range of zeros, no
  statement about the Riemann hypothesis, and no claim to have settled any open
  question of~\cite{ACLG}.
\end{enumerate}

\subsection{Revision history}
\label{sec:history}

Versions $1.0.x$ through $1.2.0$ were internal drafts, never deposited. Version
$2.0.0$ is the first archived release and the first with a different scope: two
windows rather than one. The stages are kept distinguishable on purpose --- the
versions that carried the \emph{original omission}, the version that made the
\emph{interim correction}, and the release that \emph{excludes} the affected
route. No version is rewritten in place.

\begin{itemize}
\item[\textbf{status}] Versions $1.0.x$--$1.2.0$ are \textbf{internal drafts};
  none was deposited, released or archived. Version $2.0.0$ is the first
  archived release.
\item[\textbf{1.0.x}] \textbf{carried the omission.} The upper bound
  $1.9014\times10^{-17}$ and everything derived from it rested on an incomplete
  error budget. Those versions should not be cited for the upper bound.
\item[\textbf{1.1.0}] \textbf{interim correction.} The affected claims were
  withdrawn and the endpoint reverted to a quadrature-free route.
\item[\textbf{1.2.0}] \textbf{improved, repair still open.} The
  quadrature-free endpoint was tightened past the withdrawn value. The outer
  quadrature remainder itself is \emph{not} repaired; that would be a later
  version.
\end{itemize}

\begin{itemize}
\item[\textbf{1.0.0}] (2026-11-10) First version.
\item[\textbf{1.0.1}] (2026-11-10) Citation and provenance revision only. Adds
  the two earlier technical notes whose implementation and parity conventions
  this computation inherits (\cite{Sensitivity} and~\cite{ParityNote}), the
  lineage paragraph in Section~\ref{sec:prior}, and the corresponding citations
  in Sections~\ref{sec:setup} and~\ref{sec:upper}. No numerical result, bound,
  artifact or limitation changed; the certified bracket and the parity
  statement are exactly those of version 1.0.0.
\item[\textbf{1.0.2}] (2026-11-10) Corrections after reading~\cite{Zhu} in full,
  which had previously been consulted only through its abstract and rendered
  HTML. Three attributions were too weak and are strengthened: \cite[Thm.~6.2]{Zhu}
  certifies the parity \emph{separation} at this very window, with a far stronger
  odd-sector constant, and establishes simplicity and evenness of the ground
  state (Sections~\ref{sec:prior} and~\ref{sec:parity});
  \cite[\S6.2]{Zhu} already evaluates its variational upper bounds by the
  time-domain route used in Section~\ref{sec:upper}; and the supplementary
  materials of~\cite[\S17]{Zhu} are described there, rather than absent.
  Section~\ref{sec:compared} is new and compares the two sets of bounds at this
  window. Again no numerical result, bound or artifact changed.
\item[\textbf{1.1.0}] (2026-11-10) \textbf{The certified upper endpoint
  changes.} The time-domain evaluator was found to omit the truncation error of
  its outer quadrature (Section~\ref{sec:status}), so every value it produced is
  withdrawn as a certificate and the upper bound reverts to the closed-form
  witness of Section~\ref{sec:upperclosed}: $2.1130\times10^{-17}$ in place of
  $1.9014\times10^{-17}$. The bracket widens from $4.914\times10^{-18}$ to
  $7.03\times10^{-18}$, the parity separation is re-established on the
  closed-form witness with margin $1.87\times10^{-18}$, and
  Section~\ref{sec:compared} is updated. The lower bound, the odd-sector bound,
  the $R_T$ values and the ceilings are unaffected. Four further corrections:
  the Hermitian pole term for complex $f$ (Section~\ref{sec:setup}), the degree
  of the autocorrelation, outward rounding of the ceilings, and the scope of the
  verification claims (Section~\ref{sec:repro}).
\item[\textbf{1.2.0}] (2026-11-10) \textbf{The certified upper endpoint
  improves past the withdrawn one.} Rather than repair the unbounded quadrature
  remainder, the endpoint is tightened along the route that has no quadrature to
  bound: the closed-form evaluation is pushed from $28$ even cosine modes to
  $160$, giving $U = 1.6731255341\times10^{-17}$ in place of
  $2.1130\times10^{-17}$, and in place of the withdrawn time-domain value
  $1.9014\times10^{-17}$ as well. The bracket narrows to $2.64\times10^{-18}$,
  $80.9\%$ inside the published one, and the parity margin grows to
  $6.27\times10^{-18}$. The lower bound, the odd-sector bound, the $R_T$ values
  and the ceilings are unchanged. The outer quadrature remainder of
  Section~\ref{sec:status} remains open; it is no longer on the critical path.
\item[\textbf{2.0.0}] (2026-11-12) \textbf{First archived release; the scope
  becomes two windows.} New: the second window $\widetilde L$ with its own
  two-sided bracket (Section~\ref{sec:second}); the ratio
  Corollary~\ref{cor:ratio}; the monotonicity and coverage statement
  (Proposition~\ref{prop:mono}); per-sector conditioned coercivity
  (Section~\ref{sec:cond}); the even-sector count
  (Section~\ref{sec:count}); and the proved normalization dictionary to
  \cite{ACLG} (Section~\ref{sec:dict}).

  Sharpened, with no reassembly --- the frozen $T = 400$ blocks admitted larger
  Rump shifts than had been extracted: the $L = 4/5$ lower endpoint from
  $1.41\times10^{-17}$ to $1.43300048827856\times10^{-17}$, and the odd-sector
  constant from $2.3\times10^{-17}$ to $1.38569999999861\times10^{-14}$, which
  now exceeds~\cite[Thm.~6.2]{Zhu}'s odd constant by a factor $1.68$. The
  $T = 400$ even route is thereby exhausted: the certified floor and the
  independently measured ceiling $r_{400}(f_{400})$ agree to a relative
  $5.0\times10^{-5}$, so the $T = 800$ headroom is $+16.4\%$, not $+18.3\%$.

  Corrected in the exposition: published $\mu$ values are now short exact
  decimals with explicit Schur margins, because a $\mu$ bisected to the
  supremum $a - \kappa^2/d$ cannot be serialized (its margin is $\sim10^{-41}$
  of $a$) and the earlier serialization could not be replayed; every printed
  bound is now directed, replacing in particular an invalid short upper form
  $1.6731\times10^{-17}$ by $1.674\times10^{-17}$; the eigenvalue counts are
  stated as ``exactly four even'' (both sides certified) and ``at most three
  odd'' rather than as counts read off a compression of $R_T$; and the
  rung-ratio value $0.579867$ is labelled a numerical reference rather than a
  prediction of~\cite[Hyp.~3.4]{ACLG}. The front-page draft notice is replaced
  by a scope notice; the time-domain route is excluded rather than carried.
  Appendix~\ref{app:errors} records the new failure modes.
\end{itemize}

\appendix

\section{Error history}
\label{app:errors}

Errors found during this work are recorded because several of them produced
plausible wrong numbers rather than failures, and because the countermeasures
are part of what the artifacts are for. Items 1--5 are library- or
discipline-level defects; items 6--8 are corrections to claims. A longer ledger,
including retracted statements $R1$--$R13$, is kept with the artifacts.

\begin{enumerate}
\item \textbf{Pole-coordinate phase.} The form is assembled in the coordinates
  $y_n = \sigma_n x_n$ with $\sigma_n = (-1)^{\lfloor n/2\rfloor}$, so the pole
  vector requires the same phase. Omitting it is invisible at degree $0$ and
  wrong from degree $2$ onward. It produced a genuinely negative pivot
  ($-0.3949$) that would have been reported as indefiniteness. Caught by
  comparing $|F(t)|^2$ and $F(i/2)$ against direct integration; the degree-$2$
  checks of Section~\ref{sec:upper} now cover it.
\item \textbf{Digamma on wide balls.} The library digamma returns NaN on balls
  whose real part has radius $\gtrsim 0.1$, even strictly inside the right
  half-plane. A comparison loop silently skipped those cells and returned
  $M \approx 1831$ instead of $8.83$ --- and the wrong value was used. Replaced
  by the explicit recurrence $\psi(z) = \psi(z+K) - \sum_{k<K}(z+k)^{-1}$ with
  adaptive $K$, and by a covering-free monotonicity bound for $M$.
\item \textbf{Bessel functions on wide balls.} Interval evaluation of $j_n$ on
  the complex ellipses is useless here (width $\sim10^{32}$ at $|w|\sim80$,
  unchanged at $3000$ bits). Replaced by the uniform
  Cauchy--Schwarz bound~\eqref{eq:cs}.
\item \textbf{Interval factorization.} Naive interval $LDL^{\mathsf T}$ is
  unusable at these dimensions (radii amplify about $6.5\times$ per pivot,
  reaching $O(1)$ by step $47$). Replaced by the Rump preconditioned congruence.
\item \textbf{Precision ownership.} The Arb context precision is a process
  global. Routines set it on entry and did not restore it, so a callee silently
  changed its caller's context; this produced four separate wrong constants in
  this project, each plausible-looking rather than an error (a complement
  constant at $[\pm 7.4\times10^{132}]$, an even-sector $d$ at $[\pm
  2.8\times10^{183}]$, and two coefficient sets built at $53$ bits). The repair
  is an owned-scope wrapper that asserts restoration on normal return
  \emph{and} on exception, caches keyed by precision so a low-precision value
  can never be served to a high-precision caller, and an explicit
  \emph{unresolved} result for any wide or non-finite enclosure. The lesson, and
  the harder half, is that \emph{an input constructed outside the scope carries
  the ambient precision's radius, and entering a higher-precision scope
  afterwards does not recover it}: own your inputs as well as your scope.
\item \textbf{The upper endpoint, twice mis-quoted.} Two distinct errors, both
  in the safe direction numerically but both wrong as stated. (a) Internal
  drafts quoted the enclosure \emph{midpoint}
  $1.90138291114546440\times10^{-17}$ as the bracket's upper endpoint; a
  midpoint is not an upper bound for its own ball. (b) A separate record quoted
  $1.9013829111454646892\times10^{-17}$ as ``the exact upper endpoint''; that
  value came from re-parsing the serialized ball at the ambient $53$ bits,
  which rounds the midpoint and then inflates the endpoint by about
  $2.8\times10^{-33}$ --- the precision-ownership defect of item~5 recurring in
  the very routine that documents it. It recurred twice more while that was
  being repaired: once extracting an endpoint after an owned scope had exited,
  once parsing a stored ball string at the ambient precision. Four occurrences,
  which is why deserialisation and endpoint extraction now have their own entry
  points in the precision module instead of being written inline.
  See Remark~\ref{rem:endpoint}.
\item \textbf{An omitted quadrature remainder.} The time-domain evaluator adds
  an explicit bound for the integrable head on $[0,\eta]$ but none for the
  truncation error of its outer panelled Gauss--Legendre rule, while its module
  docstring asserted that a Bernstein-ellipse remainder was used. Ball
  arithmetic does not supply that term: it encloses the nodes, weights and
  rounding, not the difference between the rule and the integral. The published
  radii were therefore incomplete, and the bound they supported is withdrawn
  (Section~\ref{sec:status}). The certified endpoint now comes from a route that
  performs no quadrature. This is the one item in this list that changed a
  headline number.
\item \textbf{A float comparison asserting an unreachable tolerance.} An
  auxiliary verifier check compared the Arb assembly with the repository's
  \src{mpmath} assembly by converting both to ordinary doubles and requiring
  agreement below $10^{-30}$. Doubles carry about sixteen digits, so the
  comparison could only report exact equality or a difference near $10^{-16}$;
  the threshold was unreachable and the check was weaker than it read. It now
  compares at full precision against the tolerance its inputs support.
\item \textbf{A retracted sparsity target.} An earlier proposed next step --- a
  power-saving sparsity bound on the near-resonant set of the prime comb --- is
  false, and is retracted in full: by continuous Kronecker--Weyl equidistribution
  that set has positive density, confirmed numerically here to four significant
  figures across four decades of $T$. The replacement is the empty-exterior gate
  of Section~\ref{sec:lower}, which needs no such estimate.
\item \textbf{Overstated interpretations, corrected.} Four further claims made
  alongside earlier versions were withdrawn or corrected: that no bound on the
  measure of a bad set could control $\int_E|F|^2$ (false --- standard
  finite-measure and local-occupation estimates do exactly that); that a Schur
  \emph{threshold} was a lower bound (it is a break-even point, and the final
  test needs strict inequalities, hence~\eqref{eq:schur}); that lowering $\beta$
  would buy a larger $a$ (refuted by Proposition~\ref{prop:beta} --- the
  $\beta\|f\|_2^2$ term had been ignored); and an accusation that the library's
  adaptive integrator returned a silently wrong value (withdrawn --- it could not
  be reproduced; the wrong numbers came from a run that still carried the
  pole-phase error of item~1, at a different cutoff, with output from two runs
  conflated).
\item \textbf{Precision ownership, in new drivers (version 2.0.0).} The defect
  of item~9 recurred four times in driver code written outside the repaired
  \src{constants()} path, each time producing a silently vacuous or
  unreplayable record rather than an error. (a) Certificate constants
  recomputed inline, before any routine had set the working precision, gave
  $d$ a radius of $3\times10^{+230}$ and $\kappa = \mathrm{NaN}$ in one sector
  only --- call-order dependent --- so the Schur step could not close and the
  sector reported a vacuous $\mu \ge 0$. (b) A one-sided shift search that
  doubled upward and stopped at the first failure certified \emph{nothing} when
  the start guess was $1.4\%$ above the true $\lambda_{\min}$. (c) Serialized
  matrix entries re-parsed after the owned scope had exited carried the ambient
  $53$-bit radius: entry radii inflated from $5.12\times10^{-49}$ to
  $7.18\times10^{-19}$, a factor $1.4\times10^{30}$, and no shift could be
  certified at all. (d) A $\mu$ bisected to the Schur supremum has margin
  $a - \mu^\ast = \kappa^2/d \approx 10^{-58}$, i.e.\ $\sim10^{-41}$ of $a$;
  serialized at $25$ digits its truncation radius $\sim3\times10^{-46}$ is
  $10^{12}$ times the margin, so the published $\mu$ was a ball straddling $a$
  and the strict test $a > \mu$ \emph{failed} on replay.

  All four were failures \emph{to certify}, not false certifications: the
  positive-definiteness routine sets its own precision, so every shift it
  reported as passing was valid, and in (d) the underlying certification had
  been performed correctly at full precision while only its serialization was
  unusable. The countermeasures are now rules: every driver takes its constants
  through \src{constants()} under owned precision with tightness assertions;
  every re-parse of a serialized ball names its precision; shift searches are
  two-sided; and a published constant must carry a margin large enough to
  survive serialization, which is why the $\mu$ values of this version are
  short exact decimals at $a(1-10^{-12})$ rather than at the supremum. This is
  the only blanket statement made here: it covers the four items above and is
  not a guarantee about defects not yet found.
\item \textbf{A downward-rounded upper bound (version 2.0.0).} The short form
  $1.6731\times10^{-17}$ was printed in several places as an upper bound for
  $U_0 = 1.67312553407068\ldots\times10^{-17}$, which it is not: a decimal
  rounded \emph{down} cannot bound a quantity from above. This is item~6
  recurring. Corrected to $1.674\times10^{-17}$, the value the artifact already
  carried, and the release gate now checks the direction of every short bound
  and explicitly rejects both $1.6731\times10^{-17}$ as an upper form and
  $0.6176$ as a short lower form for $0.617553710936882$.
\item \textbf{Counts asserted from the wrong operator (version 2.0.0).} The
  statement that exactly four even and exactly three odd eigenvalues lie below
  $10^{-2}$ was first made from the eigenvalues of a compression of the
  comparison form $R_T$. Since $Q \ge R_T$, those bound the eigenvalues of $Q$
  from neither side for $j \ge 2$, so the claim was unsupported. The even count
  is now certified on both sides (Theorem~\ref{thm:exactly4}); in the odd sector
  only ``at most three'' is claimed, and the compression counts are labelled
  numerical.
\end{enumerate}

\begin{thebibliography}{99}

\bibitem{ACLG} R. Andrews, A. Connes and A. La Gatta,
  \emph{The uncertainty operator, the Weil ground state and the zeros of
  $\zeta$: a numerical study and conditional positivity of the Weil quadratic
  form}, arXiv:2611.13129v1 (2026).
  \url{https://arxiv.org/abs/2611.13129v1}

\bibitem{Bombieri} E. Bombieri,
\emph{Problems of the Millennium: the Riemann Hypothesis},
Clay Mathematics Institute, 2000.

\bibitem{CC21} A. Connes and C. Consani,
\emph{Weil positivity and trace formula: the archimedean place},
Selecta Mathematica \textbf{27} (2021), no.~4, art.~77.

\bibitem{CCM} A. Connes, C. Consani and H. Moscovici,
\emph{Zeta zeros and prolate wave operators},
preprint, arXiv:2203.02261.

\bibitem{CvS} A. Connes and W. D. van Suijlekom,
\emph{Spectral truncations in noncommutative geometry and operator systems},
Communications in Mathematical Physics \textbf{383} (2021), 2021--2067.

\bibitem{Rump} S. M. Rump,
\emph{Verification methods: rigorous results using floating-point arithmetic},
Acta Numerica \textbf{19} (2010), 287--449.

\bibitem{Arb} F. Johansson,
\emph{Arb: efficient arbitrary-precision midpoint-radius interval arithmetic},
IEEE Transactions on Computers \textbf{66} (2017), no.~8, 1281--1292.

\bibitem{Slepian} D. Slepian and H. O. Pollak,
\emph{Prolate spheroidal wave functions, Fourier analysis and uncertainty I},
Bell System Technical Journal \textbf{40} (1961), 43--63.

\bibitem{Weil} A. Weil,
\emph{Sur les ``formules explicites'' de la th\'eorie des nombres premiers},
Comm. S\'em. Math. Univ. Lund (1952), 252--265.

\bibitem{Zhu} X. Zhu,
\emph{Weil positivity in compact windows: a finite reduction, certified
two-sided bounds, and a Landau--Widom decay law},
preprint, arXiv:2608.24827v2 (2026).

\bibitem{Sensitivity} E. De Jes\'us,
\emph{A $\delta^2$--$L^3$ Sensitivity Law for Finite-Window Weil Positivity,
with a Function-Field Control of the Connes--van Suijlekom Pipeline},
technical note, Zenodo, 2026.
\href{https://doi.org/10.5281/zenodo.23250932}{doi:10.5281/zenodo.23250932}

\bibitem{ParityNote} E. De Jes\'us,
\emph{Parity Blindness of Finite-Window Weil Detectors at a Real Zero:
the Even Sector Cannot See a Siegel Configuration and the Odd Sector Is
Maximally Sensitive to It},
technical note, version 2, Zenodo, 2026.
\href{https://doi.org/10.5281/zenodo.23624161}{doi:10.5281/zenodo.23624161}

\bibitem{Repo} E. De Jes\'us,
\emph{weil-window-control}, research repository, branch
\texttt{positive-weil-certificate-2026-11-09}.\\
\url{https://github.com/innerlightr-wq/weil-window-control}

\end{thebibliography}

\end{document}
